\chapter{Guards and public languages}
\label{prov:ch:guard}
\label{guard:ch}

The native scalar core has two introduction rules: \rulename{DI} by cases,
at every matrix, and \rulename{QI} by progress, at literal equations
(\cref{base:sec:kernel}).  This chapter treats calculi with a single
introduction of $\D{}{}$.  Its premise retains one local assumption, the
instance of the matrix at zero (\emph{cases}) or the instance at the
eigenvariable (\emph{progress}), and the choice between the two is made by
the shape of the matrix alone: no certificate bit, no test of the witness.
At a literal equation the two modes prove different formulas.  They prove
the same formula when the instance at zero is provable, and a \emph{guard}
is a literal shape of an equation whose instance at zero is provable by
computation and into which every equation can be lifted.  The guard is used
in two ways.  Either the formula language is restricted to a \emph{public
language} in which every atomic matrix below $\D{}{}$ has the guarded shape,
and an arbitrary native formula is translated into it; or the full language
is kept and the progress mode is used exactly at the guarded shape.

The chapter is organized by idea.  \Cref{guard:sec:shapes} defines the
guarded calculi, the shapes, and what substitution does to them.
\Cref{guard:sec:bootstrap} recovers numerical cases, implication between
equations and equation induction from a guard, once for all guards, and
translates the native language into a public language with both composites.
\Cref{guard:sec:sources} determines which transport sources are needed: the
literal $0$ is not.  \Cref{guard:sec:composition} shows that the two
restrictions together are incomplete in a strictly guarded language, and
gives four repairs: auxiliary formulas, the administrative orbit, the mirror
guard, and protected atoms.  \Cref{guard:sec:full} treats the guards that
work on the full native language, the diagonal guard and the mirror guard
with six rules.  \Cref{guard:sec:compare} compares the guards, and
\cref{guard:sec:checkers} records the checkers and their sizes.

Throughout the chapter $1$ abbreviates $\Suc 0$ and $2$ abbreviates
$\Suc\Suc 0$; $s\equiv t$ is literal identity of terms; a term is
\emph{iterator-headed} if it has the form $\It{n}{a}{u}{h}$ and
\emph{successor-headed} if it has the form $\Suc t$; and
\[
 W(r)\defeq\It{0}{r}{u}{0}
\]
is the \emph{wrapper} of $r$, an iterator-headed term with $W(r)=r$ by the
first computation equation.  For a calculus $\mathsf K$ and formulas $X$,
$Y$ of its language, $X\dashv\vdash_{\mathsf K}Y$ means that
$\Gamma\proves_{\mathsf K}X$ if and only if $\Gamma\proves_{\mathsf K}Y$,
for every context $\Gamma$ of $\mathsf K$; the subscript is omitted when the
calculus is clear.

In \crefrange{guard:sec:shapes}{guard:sec:composition} transport is the
uncrossed rule in a formula template,
\[
 \infrule{\Gamma\proves s=t\qquad\Gamma\proves F(s)}{\Gamma\proves F(t)},
\]
as the packages of those sections state it.  For $\Nat$ it is a derived
rule (\cref{base:cor:tr}).  It is the rule written $\rulename{Tr}_F$ in
\cref{kern:sec:dispatch}, and its restriction to the sources
$s\not\equiv 0$, written $\rulename{Tr}^{-}$ below, is the rule
$\rulename{Tr}_F^{-}$ of that section, not the crossed rule of
\cref{kern:def:OM}.  The orientation of the
conclusion matters as soon as the sources $s$ are restricted
(\cref{guard:rem:crossed}).  In \cref{guard:sec:full} the equality rule is
the rule \rulename{Tr} of the kernel.

% ---------------------------------------------------------------------
\section{Guards as literal shapes}
\label{guard:sec:shapes}

This section collects the definitions common to \pkg{guarded\_core},
\pkg{guarded\_source}, \pkg{guarded\_orbit},
\pkg{deterministic\_auxiliary}, \pkg{mirror\_guard}, \pkg{mirror\_core},
\pkg{protected\_core} and \pkg{diagonal\_guard}.

\subsection{The root-determined introduction}

\begin{definition}[Root-determined introduction]
\label{guard:def:root}
For a matrix $B(x,y)$ put
\[
 \eta_B(z)\defeq
 \begin{cases}
  z, & B \text{ is a literal equation},\\
  0, & B \text{ is compound (its root is } \Q{} \text{ or } \D{}{}).
 \end{cases}
\]
The \emph{root-determined introduction} is the rule
\[
 \infrule{\Gamma,\ B(t,\eta_B(z))\proves B(t,\Suc z)}
         {\Gamma\proves\D{x}{y}B(x,y)}\ \rulename{DI}_{\mathrm{root}},
\]
where $z$ is not free in $\Gamma$, in the witness $t$ or in $\D{x}{y}B$, and
the displayed assumption is discharged.  The rule is in \emph{progress
mode} at a literal equation and in \emph{case mode} at a compound matrix.
\end{definition}

The mode is read off the matrix before the witness is substituted, and it is
the same for every substitution instance of the matrix, because substitution
of terms for variables preserves and reflects the root of a formula.

\begin{definition}[Guarded calculus]
\label{guard:def:calculus}
Let $\mathcal F$ be a set of native formulas, the \emph{proof language}, and
$\mathcal T\subseteq\mathcal F$ a set of \emph{templates}.  Both contain
every equation and are closed under substitution of terms for free
variables; $\mathcal F$ contains $B(t,r)$ and $\Q{y}B(t,y)$ whenever it
contains $\D{x}{y}B(x,y)$, and $B(t)$ whenever it contains $\Q{y}B(y)$.  Let
$\mathcal S$ be a set of terms closed under injective renaming of variables,
the \emph{admissible sources}.  The \emph{guarded calculus}
$\mathsf K(\mathcal F,\mathcal S)$ derives sequents $\Gamma\proves A$, with
$\Gamma$ a finite list of formulas of $\mathcal F$ and $A\in\mathcal F$, by
the following seven rules.
\begin{gather*}
 \infrule{A\in\Gamma}{\Gamma\proves A}\ \rulename{Hyp}
 \qquad
 \infrule{}{\Gamma\proves\It{0}{a}{u}{h}=a}\ \rulename{I0}
 \qquad
 \infrule{}{\Gamma\proves\It{\Suc n}{a}{u}{h}=h\subst{\It{n}{a}{u}{h}}{u}}\
 \rulename{IS}
 \\[1.2ex]
 \infrule{\Gamma\proves s=t\qquad\Gamma\proves F(s)}
         {\Gamma\proves F(t)}\ \rulename{Tr}_{\mathcal S}
 \quad (s\in\mathcal S,\ F\in\mathcal T)
 \qquad
 \infrule{\Gamma\proves\Q{y}B(y)\qquad\Gamma\proves B(0)}
         {\Gamma\proves B(t)}\ \rulename{QE}
 \\[1.2ex]
 \infrule{\Gamma\proves\D{x}{y}B(x,y)\qquad
          \Gamma,\ \Q{y}B(b,y)\proves A}
         {\Gamma\proves A}\ \rulename{DE}
 \qquad
 \infrule{\Gamma,\ B(t,\eta_B(z))\proves B(t,\Suc z)}
         {\Gamma\proves\D{x}{y}B(x,y)}\ \rulename{DI}_{\mathrm{root}}
\end{gather*}
In \rulename{DE} the variable $b$ is not free in $\Gamma$, $A$ or
$\D{x}{y}B$; in $\rulename{DI}_{\mathrm{root}}$ the conclusion belongs to
$\mathcal F$ and $z$ is not free in $\Gamma$, $t$ or $\D{x}{y}B$.  There is
no reflexivity axiom, no successor axiom, no rule of induction, no rule that
introduces $\Q{}$, and no structural rule.
\end{definition}

Three classes of sources occur: all terms; the terms $s\not\equiv 0$; and
the iterator-headed terms.  The condition $s\not\equiv 0$ is a test of
syntax: a variable, or an iterator whose value is zero, is an admissible
source.  The rule with the sources $s\not\equiv 0$ is written
$\rulename{Tr}^{-}$ and the rule with iterator-headed sources
$\rulename{Tr}_I$.  Two kinds of proof language occur: the full native
language $\mathcal L$, with $\mathcal T=\mathcal L$; and the proof language
$\mathcal P^{Q}$ of a public language $\mathcal P$
(\cref{guard:def:public}), in which $\Q{}$ occurs only at the root of a
hypothesis and the templates are the equations and the formulas of
$\mathcal P$.

\begin{remark}[Source-selected elimination]
\label{guard:rem:se}
With the sources $s\not\equiv 0$ the two rules $\rulename{Tr}^{-}$ and
\rulename{QE} are the two branches of the rule \rulename{SE} of
\cref{kern:sec:dispatch}, whose first premise is $\Q{v}F(v)$ at the source
$0$ and $s=t$ at every other source.  In this form, which is the one of
\pkg{guarded\_source}, \pkg{guarded\_orbit} and
\pkg{deterministic\_auxiliary}, a guarded calculus has six rule patterns:
the source selects the elimination and the root of the matrix selects the
introduction.  The calculus $\mathsf Z$ of \cref{kern:def:Z} has
\rulename{SE} with the two introductions of $\Nat$; the calculi of this
chapter have it with $\rulename{DI}_{\mathrm{root}}$ alone.
\end{remark}

\begin{remark}[Clause contexts]
\label{guard:rem:clauses}
\pkg{guarded\_core} and \pkg{mirror\_guard} state their calculi with the
clause contexts of \cref{base:sec:clauses}: \rulename{DE} adds the clause
$[B(b,0);\,y.\,B(b,y)]$, an assumption is a constant clause, and one rule
\rulename{Use} infers $B(t)$ from a derivation of $A$ and a clause
$[A;\,y.\,B(y)]$ of the context.  A hypothesis $\Q{y}B(y)$ used by
\rulename{Hyp} and \rulename{QE} is the clause $[B(0);\,y.\,B(y)]$ used by
\rulename{Use}, so the calculi below are the clause calculi of those
packages for contexts of assumptions and opened clauses.  The statements and
proofs of this chapter hold verbatim for contexts with arbitrary clauses
$[A;\,y.\,B(y)]$ of public formulas, with a translation of clauses
componentwise.
\end{remark}

\subsection{Structural operations, the kernel, and equality at admissible
sources}

\begin{lemma}[Structural operations]
\label{guard:lem:struct}
In every guarded calculus $\mathsf K(\mathcal F,\mathcal S)$ weakening,
injective renaming of free variables and cut are admissible: if
$\Gamma\proves A$ and $\Delta\proves C$ for every $C\in\Gamma$, then
$\Delta\proves A$.  If $\mathcal S$ is closed under substitution, as are the
class of all terms and the class of iterator-headed terms, then substitution
is admissible: $\Gamma\proves A$ implies
$\Gamma[\theta]\proves A[\theta]$.
\end{lemma}

\begin{proof}
Each statement is proved by induction on the derivation, with the
eigenvariables of \rulename{DE} and $\rulename{DI}_{\mathrm{root}}$ renamed
away from the data that are inserted, as for $\Nat$
(\cref{base:sec:structural}).  Weakening and cut do not change the matrix of
an introduction or the source of a transport, so every inference keeps its
rule and its mode.  A renaming, and a substitution $\theta$, preserve the
root of every formula, so an introduction in progress mode remains one in
progress mode and likewise for case mode; the languages $\mathcal F$ and
$\mathcal T$ are closed under substitution; and the source $s[\theta]$ of a
transported equation is admissible when $\mathcal S$ is closed under
$\theta$, which holds for every injective renaming and, in the two named
cases, for every substitution.
\end{proof}

The class of terms $s\not\equiv 0$ is not closed under substitution: a
variable is an admissible source and its instance $0$ is not.  A derivation
with $\rulename{Tr}^{-}$ therefore cannot be substituted into rule by rule.
The arguments of \cref{guard:sec:sources,guard:sec:composition} are arranged
so that no such substitution is used before the unrestricted rule has been
recovered.

\begin{lemma}[Inclusion in the kernel]
\label{guard:lem:incl}
Every derivation of a guarded calculus $\mathsf K(\mathcal F,\mathcal S)$
is transformed, rule by rule, into a derivation of the same sequent in
$\Nat$.
\end{lemma}

\begin{proof}
\rulename{Hyp}, \rulename{I0}, \rulename{IS}, \rulename{QE} and \rulename{DE}
are rules of $\Nat$.  $\rulename{Tr}_{\mathcal S}$ is an instance of
transport in a formula template (\cref{base:cor:tr}).
$\rulename{DI}_{\mathrm{root}}$ at a compound matrix is \rulename{DI}, and
at a literal equation it is the witnessed introduction by progress, which
$\Nat$ derives (\cref{base:cor:combined}(i)).
\end{proof}

Consequently every guarded calculus is sound for the ordinary reading
(\cref{base:sec:reading}): a derivable sequent of
$\mathsf K(\mathcal F,\mathcal S)$ is derivable in $\Nat$, hence its reading
is provable in $\Teq$ and in $\iPRA$ (\cref{base:thm:bridge}).

\begin{lemma}[Equality at admissible sources]
\label{guard:lem:safe}
Let $\mathsf K=\mathsf K(\mathcal F,\mathcal S)$ and let $\mathcal S$ contain
every iterator-headed term.  The following hold over every context $\Gamma$,
by derivations that use only \rulename{I0}, \rulename{IS} and
$\rulename{Tr}_{\mathcal S}$ in equation templates.
\begin{enumerate}[label=\textup{(\roman*)}]
\item $\Gamma\proves r=r$ for every term $r$.
\item If $\Gamma\proves p=q$ and $p\in\mathcal S$, then $\Gamma\proves q=p$;
 $\Gamma\proves f(p)=f(q)$ for every term template $f(v)$;
 $\Gamma\proves r=p$ implies $\Gamma\proves r=q$; and $\Gamma\proves E(p)$
 implies $\Gamma\proves E(q)$ for every equation template $E(v)$.
\item Every instance of \rulename{I0} and \rulename{IS} can be reversed.
\item \emph{Join at zero.}  If $p,q\in\mathcal S$, $\Gamma\proves p=0$ and
 $\Gamma\proves q=0$, then $\Gamma\proves p=q$.
\item \emph{Lifting.}  Let $p,a,b\in\mathcal S$, let $q\in\mathcal S$ or
 $q\equiv 0$, and let $\Gamma\proves a=p$ and $\Gamma\proves b=q$.  Then
 $\Gamma\proves p=q$ if and only if $\Gamma\proves a=b$.  The implication
 from $a=b$ to $p=q$ needs no condition on $q$.
\item \emph{Expansion on the right.}  Let $\mathcal S$ contain every term
 other than the literal $0$, let $q\not\equiv 0$ and $\Gamma\proves q=r$.
 Then, for every term $s$, $\Gamma\proves s=q$ if and only if
 $\Gamma\proves s=r$.
\end{enumerate}
\end{lemma}

\begin{proof}
(i)--(iv) For the sources $s\not\equiv 0$ these are
\cref{kern:lem:zsafe,kern:lem:zjoin}; the wrapper $W(r)$ is the term $J(r)$
used there.  The derivations given there transport only at the sources $p$
and $q$ named in each item and at $W(0)$, $W(p)$, $W(q)$ and $W(r)$, and
the left side of a computation equation is iterator-headed.  Hence they are
derivations of $\mathsf K$ for every class $\mathcal S$ that contains the
iterator-headed terms, under the stated conditions $p,q\in\mathcal S$.  No
equation with the source $0$ is used, and in particular transitivity
through the literal $0$ is avoided.

(v) Let $\Gamma\proves p=q$.  Transitivity through $p$ gives $a=q$.  If
$q\in\mathcal S$, reverse $b=q$ at the source $b$ and pass through $q$:
$a=b$.  If $q\equiv 0$, then $a=0$ and $b=0$, and (iv) gives $a=b$.
Conversely let $\Gamma\proves a=b$.  Transport along $a=p$ in the template
$v=b$ gives $p=b$, and transport along $b=q$ in the template $p=v$ gives
$p=q$; the sources are $a$ and $b$.

(vi) From $s=q$, transport along $q=r$ in the template $s=v$ gives $s=r$.
Conversely let $\Gamma\proves s=r$.  If $r\not\equiv 0$, reverse $q=r$ at
the source $q$ and transport at the source $r$ in the template $s=v$.  If
$r\equiv 0$ and $s\not\equiv 0$, then $s=0$ and $q=0$, and (iv) gives $s=q$.
If $r\equiv 0$ and $s\equiv 0$, the required equation $0=q$ is the reverse
of $q=0$ at the source $q$.  The case distinction is on the syntax of the
terms $r$ and $s$; it is not a case distinction inside the calculus.
\end{proof}

With all terms as sources, (i) and (ii) are the ordinary equality rules,
and the calculus has reflexivity, symmetry, transitivity, congruence and
transport in every equation template.

\subsection{Deleting the choice is not conservative}

The administrative calculus is
$\mathsf U\defeq\mathsf K(\mathcal L,\text{all terms})$: the full native
language with the root-determined introduction.  It is the calculus that the
compact checkers of \pkg{guarded\_core} implement, and by
\cref{guard:lem:incl} it is a subcalculus of $\Nat$.  It is a proper one.

\begin{definition}[Mixed reading]
\label{guard:def:mixed}
The \emph{mixed reading} $A^\dagger$ of a native formula is
\begin{align*}
 (s=t)^\dagger &\defeq (s=t),\\
 (\Q{y}B)^\dagger &\defeq B^\dagger(0)\to\forall y\,B^\dagger(y),\\
 (\D{x}{y}B)^\dagger &\defeq
 \begin{cases}
  \exists x\,\forall y\,\bigl(B(x,y)\to B(x,\Suc y)\bigr),
    & B \text{ a literal equation},\\
  \exists x\,\bigl(B^\dagger(x,0)\to\forall y\,B^\dagger(x,y)\bigr),
    & B \text{ compound}.
 \end{cases}
\end{align*}
\end{definition}

The mixed reading depends on the syntax of the matrix, not on its meaning:
it is not invariant under the replacement of a matrix by an interderivable
one.  It commutes with substitution, because substitution preserves the
root of a matrix.

\begin{lemma}[Soundness for the mixed reading]
\label{guard:lem:mixed-sound}
If $\Gamma\proves_{\mathsf U}A$, then $\Teq$ proves $A^\dagger$ from
$\Gamma^\dagger$.
\end{lemma}

\begin{proof}
By induction on the derivation.  \rulename{Hyp}, the computation equations
and \rulename{QE} are immediate, and transport is sound because the reading
commutes with substitution.  At an atomic $\rulename{DI}_{\mathrm{root}}$
the premise proves $B(t,z)\to B(t,\Suc z)$ for a fresh $z$, and
generalization on $z$ followed by existential introduction gives the
progressive reading.  At a compound $\rulename{DI}_{\mathrm{root}}$ the
premise proves $B^\dagger(t,0)\to B^\dagger(t,\Suc z)$, and coverage
$y=0\vee\exists z\,y=\Suc z$ extends it to
$B^\dagger(t,0)\to\forall y\,B^\dagger(t,y)$.  For \rulename{DE} at a
literal equation $E$, fix $x$ with $\forall y\,(E(x,y)\to E(x,\Suc y))$;
induction for the literal equation $E(x,y)$ in $\Teq$ gives
$E(x,0)\to\forall y\,E(x,y)$, which is the reading of the hypothesis
$\Q{y}E(b,y)$ that the rule supplies.  \rulename{DE} at a compound matrix is
existential elimination.
\end{proof}

\begin{proposition}[The root-determined introduction alone is weaker]
\label{guard:prop:strict}
With $x$ not free in the matrix, $\Nat$ derives $\D{x}{y}(y=1)$ and
$\mathsf U$ does not.
\end{proposition}

\begin{proof}
The rule $\rulename{DI}_{\mathrm{root}}$ is the introduction by cases
restricted to compound matrices together with the witnessed introduction by
progress at literal equations.  Hence $\mathsf U$ is the calculus of
\cref{kern:prop:compoundonly}(ii) with transport in a formula template, and
the statement is that proposition.  Its interpretation is the truth of the
mixed reading in the natural numbers, which \cref{guard:lem:mixed-sound}
strengthens to provability in $\Teq$: the mixed reading
$\exists x\,\forall y\,(y=1\to\Suc y=1)$ of the formula fails at $y=1$.
\end{proof}

\begin{remark}
\label{guard:rem:strict-variants}
\pkg{guarded\_core} gives the example with the matrix $p(y)=1$, where
$p(y)\defeq\It{y}{0}{u}{\sel(u,1,0)}$ has the values $0,1,0$ at $0,1,2$: the
introduction by cases has the base $0=1$, and the progressive reading fails
between $1$ and $2$.  The formula $\D{x}{y}(y=1)$ is the one of its
implementation section.  The formula of the proposition shows two readings
of one raw binder that disagree in the natural numbers, so no formula
composite can identify the atomic binder of $\mathsf U$ with its ordinary
reading; no impossibility for other interpretations of $\mathsf U$ is
claimed.
\end{remark}

Hence the single introduction is complete only together with a restriction:
of the language, so that the two readings agree at every atomic matrix; or of
the progress mode, to matrices at which they agree.

\subsection{The shapes}

\begin{definition}[Guarded shapes]
\label{guard:def:shapes}
Let $B(x,y)$ be a literal equation with the distinguished variables $x$ and
$y$.  It is
\begin{enumerate}[label=\textup{(\alph*)}]
\item \emph{zero-seeded in $y$} if it is literally
 $\It{y}{0}{u}{h}=\It{y}{0}{u'}{k}$;
\item \emph{mirrored in $y$} if it is literally $a=\It{y}{a}{u}{h}$, the
 left side and the seed being the same raw term $a$;
\item \emph{diagonal in $y$} if it is literally
 $\It{y}{y}{u}{h}=\It{y}{y}{u'}{k}$;
\item an \emph{orbit matrix in $y$} if it is literally
 $y=\It{y}{a}{u}{\Suc u}$ with $y$ not free in $a$.
\end{enumerate}
The bodies $h$ and $k$ are arbitrary: they may contain $x$, $y$, their own
state variable and further variables; in (b) the term $a$ may contain $x$
and $y$.  The two state binders of (a) and (c) are independent.  An equation
$p=q$ is \emph{protected} if $p$ and $q$ are iterator-headed.  For an
equation $E\equiv(s=t)$ and a term $v$, the \emph{gate} and the \emph{mirror
gate} of $E$ at $v$ are
\[
 G_vE\defeq\bigl(\sel(v,0,s)=\sel(v,0,t)\bigr),
 \qquad
 H_vE\defeq\bigl(s=\sel(v,s,t)\bigr),
\]
where $G_vE$ is $\operatorname{gate}(v,E)$ of \cref{base:def:gate} at an
equation; the \emph{wrapped} equation is $W(s)=W(t)$; and for an equation
$E(y)\equiv(s(y)=t(y))$ the \emph{diagonal lift} is
\[
 \Delta_yE\defeq\bigl(\It{y}{y}{u}{s(y)}=\It{y}{y}{u}{t(y)}\bigr),
 \qquad u\notin\FV(s,t).
\]
\end{definition}

For a variable $y$, the gate $G_yE$ is zero-seeded in $y$, the mirror gate
$H_yE$ is mirrored in $y$ and $\Delta_yE$ is diagonal in $y$, whether or not
$y$ occurs in $E$; in each the update does not use its state.  Each shape is
a condition on the syntax.  None asks for an arithmetical equality or for
the provability of a base case.

\begin{lemma}[Instances at zero]
\label{guard:lem:zero}
Let $\mathsf K(\mathcal F,\mathcal S)$ be a guarded calculus whose sources
include the iterator-headed terms, and let $B(x,y)$ be zero-seeded,
mirrored or diagonal in $y$.  Then $\Gamma\proves B(t,0)$ for every context
$\Gamma$ and every term $t$, by \rulename{I0} and transport in equation
templates at iterator-headed sources.
\end{lemma}

\begin{proof}
If $B$ is zero-seeded or diagonal, both sides of $B(t,0)$ have the form
$\It{0}{0}{u}{h'}$; each is equal to $0$ by \rulename{I0}, and
\cref{guard:lem:safe}(iv) joins them.  If $B$ is mirrored, $B(t,0)$ is
$a'=\It{0}{a'}{u}{h'}$, the reverse of an instance of \rulename{I0}
(\cref{guard:lem:safe}(iii)).
\end{proof}

An orbit matrix has no such property: its instance at zero is
$0=\It{0}{a}{u}{\Suc u}$, which is interderivable with $0=a$.

\begin{lemma}[Cases at a matrix with a provable base]
\label{guard:lem:atomcases}
Let $B(x,y)$ be a literal equation such that $\D{x}{y}B$ belongs to the
proof language of a guarded calculus and $\Gamma\proves B(t,0)$ for every
$\Gamma$ and $t$.  If $\Gamma,B(t,0)\proves B(t,\Suc z)$ with $z$ fresh,
then $\Gamma\proves\D{x}{y}B$.
\end{lemma}

\begin{proof}
Cut the derivation of $B(t,0)$ into the premise: $\Gamma\proves
B(t,\Suc z)$.  Weaken by $B(t,z)$ and apply $\rulename{DI}_{\mathrm{root}}$
in progress mode.  The variable $z$ stays fresh because the base that was
removed does not contain it.
\end{proof}

\begin{proposition}[Substitution and the shapes]
\label{guard:prop:subst-shapes}
Let $\sigma$ be a capture-free substitution of terms for variables other
than the distinguished $x$ and $y$, with $x$ and $y$ not free in its values.
\begin{enumerate}[label=\textup{(\roman*)}]
\item The root of a formula is preserved and reflected by $\sigma$: $B$ is a
 literal equation if and only if $B\sigma$ is.
\item The zero-seeded and the mirrored shape are preserved by $\sigma$.
 Neither is reflected: for variables $p$ and $q$ the matrix
 $\It{y}{p}{u}{0}=\It{y}{0}{u}{0}$ is not zero-seeded and becomes
 zero-seeded under $p\mapsto 0$, and the matrix $p=\It{y}{q}{u}{0}$ is not
 mirrored and becomes mirrored under $q\mapsto p$.
\item The diagonal shape is preserved and reflected: $B\sigma$ is diagonal
 in $y$ if and only if $B$ is.  The same holds for the substitution of a
 witness for $x$, if $y$ is not free in the witness.
\item The classes of protected equations and of iterator-headed terms are
 closed under $\sigma$; the class of terms other than the literal $0$ is
 not.
\end{enumerate}
\end{proposition}

\begin{proof}
(i), (ii) and (iv) are read off the definition of substitution.  In (ii) the
shapes are preserved because the literal $0$, the variable $y$ and the two
equal occurrences of $a$ are mapped to $0$, to $y$ and to two equal
occurrences of $a\sigma$.

(iii) Rename the binders away from the values of $\sigma$.  A
constructor-headed term remains headed by the same constructor.  A term
$r\sigma$ is the distinguished variable $y$ if and only if $r$ is $y$:
another variable that is kept differs from $y$; the value of a substituted
variable does not contain $y$; and a constructor-headed term does not become
a variable.  Hence an iterator $\It{n}{a}{u}{h}\sigma$ has the pair $(y,y)$
as counter and seed if and only if $\It{n}{a}{u}{h}$ has.  A side of the
equation that is a substituted variable may become iterator-headed, but its
value has no free $y$, so it does not acquire the shape.  A side that is an
iterator has the shape after $\sigma$ exactly when it has it before.  The
bodies are unrestricted and impose no condition.  The substitution of a
witness is the same argument with $x$ among the substituted variables.
\end{proof}

The proposition is the reason for the two treatments of a guard.  In a
public language the progress mode is selected by the root, which is stable,
and the shape is a requirement on the language, for which preservation
suffices.  On the full language the progress mode is selected by the shape
itself, and a shape that substitution can create changes the mode of an
inference under substitution (\cref{guard:sec:full}).

\subsection{Public languages}

\begin{definition}[Public language and its proof language]
\label{guard:def:public}
Let $\mathfrak A$ be a class of equations and $\mathfrak M$ a class of
atomic matrices, both closed under substitution, with every instance
$B(t,r)$ of a matrix $B\in\mathfrak M$ in $\mathfrak A$.  The \emph{public
language} $\mathcal P(\mathfrak A,\mathfrak M)$ is the least set of
formulas such that
\begin{enumerate}[label=\textup{(\arabic*)}]
\item every equation of $\mathfrak A$ is public;
\item $\D{x}{y}B$ is public if $B$ is a compound public formula;
\item $\D{x}{y}B$ is public if $B(x,y)$ is a matrix of $\mathfrak M$ in $y$.
\end{enumerate}
A \emph{matrix of $\mathcal P$ in $y$} is a compound public formula or a
matrix of $\mathfrak M$ in $y$.  The \emph{proof language}
$\mathcal P^{Q}$ consists of the public formulas, all equations, and the
\emph{resources} $\Q{y}B(y)$ for the matrices $B$ of $\mathcal P$ in $y$.
\end{definition}

A public language has no $\Q{}$.  Clause (3) is the restriction: an atomic
matrix immediately below $\D{}{}$ has a prescribed literal shape, and no
decision procedure and no proof is involved in recognizing it.  Compound
matrices are unrestricted within the grammar.  By
\cref{guard:prop:subst-shapes} the languages below are closed under
substitution.  The public languages of the chapter are
\begin{align*}
 \mathcal P_0 &\defeq\mathcal P(\text{all equations},\ \text{zero-seeded}),\\
 \mathcal P_0^{+} &\defeq
   \mathcal P(\text{all equations},\ \text{zero-seeded or orbit}),\\
 \mathcal P_m &\defeq\mathcal P(\text{all equations},\ \text{mirrored}),\\
 \mathcal P_W &\defeq\mathcal P(\text{protected equations},\ \text{zero-seeded}),
\end{align*}
and the calculi are listed in \cref{guard:tab:calculi}.  The designated
contradiction of $\mathcal P_0$, $\mathcal P_0^{+}$ and $\mathcal P_m$ is
$0=1$; that of $\mathcal P_W$ is $W(0)=W(1)$.

\begin{table}[ht]
\centering
\begin{tabular}{@{}llll@{}}
\toprule
Calculus & Proof language & Sources & Package\\
\midrule
$\mathsf G$ & $\mathcal P_0^{Q}$ & all terms & \pkg{guarded\_core}\\
$\mathsf U$ & $\mathcal L$ & all terms & \pkg{guarded\_core}\\
$\mathsf U^{-}$ & $\mathcal L$ & $s\not\equiv 0$ &
 \pkg{guarded\_source}, \pkg{deterministic\_auxiliary}\\
$\mathsf G^{-}$ & $\mathcal P_0^{Q}$ & $s\not\equiv 0$ &
 incomplete (\cref{guard:cor:obstruction})\\
$\mathsf G^{-}_{\mathrm{orb}}$ & $(\mathcal P_0^{+})^{Q}$ &
 $s\not\equiv 0$ & \pkg{guarded\_orbit}\\
$\mathsf M$ & $\mathcal P_m^{Q}$ & $s\not\equiv 0$ & \pkg{mirror\_guard}\\
$\mathsf H$ & $\mathcal P_W^{Q}$ & iterator-headed & \pkg{protected\_core}\\
\bottomrule
\end{tabular}
\caption{The guarded calculi $\mathsf K(\mathcal F,\mathcal S)$ of
\cref{guard:sec:shapes,guard:sec:bootstrap,guard:sec:sources,guard:sec:composition}.
Each has the seven rules of \cref{guard:def:calculus}.}
\label{guard:tab:calculi}
\end{table}

% ---------------------------------------------------------------------
\section{The guarded bootstrap and the translation into a public language}
\label{guard:sec:bootstrap}

This section draws on \pkg{guarded\_core} for the gate, the padding, the
implication between equations and the translation; on \pkg{mirror\_guard}
for the form of the induction argument that uses no substitution; on
\pkg{guarded\_orbit} and \pkg{protected\_core} for the translation of
$\Q{}$; and on \pkg{deterministic\_auxiliary} for the same translation.  The
argument is given once, for an arbitrary guard, and is instantiated in
\cref{guard:sec:core,guard:sec:composition}.

\subsection{Guard schemes}

\begin{definition}[Guard scheme]
\label{guard:def:scheme}
Let $\mathsf K=\mathsf K(\mathcal F,\mathcal S)$ be a guarded calculus and
$\mathcal P=\mathcal P(\mathfrak A,\mathfrak M)$ a public language with
$\mathcal P^{Q}\subseteq\mathcal F$.  A \emph{guard scheme} for
$\mathsf K$ and $\mathcal P$ on a class $\mathfrak E\subseteq\mathfrak A$ of
equations, closed under substitution, consists of
\begin{itemize}
\item an equation $\gamma_vU\in\mathfrak E$ for every $U\in\mathfrak E$ and
 every term $v$, and
\item a relation of \emph{compatible pairs} on $\mathfrak E$, closed under
 substitution, with an equation $\mu_w(U,V)\in\mathfrak E$ for every
 compatible pair $(U,V)$ and every term $w$,
\end{itemize}
both commuting with substitution, such that for all $U,V$, all terms $w$
and all variables $y$:
\begin{enumerate}[label=\textup{(G\arabic*)}]
\item $\gamma_yU$ is a matrix of $\mathfrak M$ in $y$, whether or not $y$
 occurs in $U$;
\item $\Gamma\proves\gamma_0U$ for every context $\Gamma$, and
 $\gamma_{\Suc w}U\dashv\vdash U$;
\item $\mu_0(U,V)\dashv\vdash U$ and $\mu_{\Suc w}(U,V)\dashv\vdash V$;
\item $(\gamma_yU,\,U)$ is a compatible pair.
\end{enumerate}
\end{definition}

Two schemes are used.  The \emph{zero-seeded scheme} has
\[
 \gamma_v(s=t)\defeq G_v(s=t)=\bigl(\sel(v,0,s)=\sel(v,0,t)\bigr),
 \qquad
 \mu_w(s=t,\ r=q)\defeq\bigl(\sel(w,s,r)=\sel(w,t,q)\bigr),
\]
with every pair compatible.  The \emph{mirror scheme} has
\[
 \gamma_v(s=t)\defeq H_v(s=t)=\bigl(s=\sel(v,s,t)\bigr),
 \qquad
 \mu_w(s=a,\ s=b)\defeq\bigl(s=\sel(w,a,b)\bigr),
\]
with $(U,V)$ compatible when $U$ and $V$ have the same left side.  In both,
(G1) and (G4) hold by inspection: $G_yU$ is zero-seeded and $H_yU$ is
mirrored in $y$, and $H_y(s=t)$ has the left side $s$.  Conditions (G2) and
(G3) are statements about a calculus and are verified where the scheme is
used.

Fix a calculus, a public language and a guard scheme on $\mathfrak E$.  For
$U\in\mathfrak E$ and a compatible pair $(U,V)$ put
\[
 L(U)\defeq\D{a}{v}\gamma_vU,
 \qquad
 (U\Rightarrow V)\defeq\D{a}{w}L\bigl(\mu_w(U,V)\bigr),
\]
with fresh bound variables, $a$ not occurring in the matrix.  Both are
compound public formulas, by (G1).  For a formula $A$ that is an equation of
$\mathfrak E$ or a compound public formula, the \emph{padding} is
\[
 P(A)\defeq
 \begin{cases}
  L(A), & A\in\mathfrak E,\\
  A, & A \text{ compound}.
 \end{cases}
\]
$P(A)$ is always compound.

\begin{lemma}[Padding]
\label{guard:lem:pad}
$L(U)\dashv\vdash U$ for every $U\in\mathfrak E$, and hence
$P(A)\dashv\vdash A$.
\end{lemma}

\begin{proof}
Let $\Gamma\proves U$.  The matrix $\gamma_vU$ is a literal equation, so the
introduction of $L(U)$ is in progress mode; with the witness $0$ its premise
is $\Gamma,\gamma_zU\proves\gamma_{\Suc z}U$, which follows from $U$ by
(G2) without the local assumption.  Conversely let $\Gamma\proves L(U)$.
\rulename{DE} adds the resource $\Q{v}\gamma_vU$; its base $\gamma_0U$ is
derivable by (G2), \rulename{QE} at $1$ gives $\gamma_{1}U$, and (G2) gives
$U$.  The eigenvariable of \rulename{DE} does not occur in $U$.
\end{proof}

\begin{lemma}[Numerical cases]
\label{guard:lem:nc}
Let $A(y)$ be a compound public formula or an equation of $\mathfrak E$, and
let $z$ be not free in $\Gamma$ and $A$.  If $\Gamma\proves A(0)$ and
$\Gamma\proves A(\Suc z)$, then $\Gamma\proves A(t)$ for every term $t$.
\end{lemma}

\begin{proof}
Let $A$ be compound.  Introduce $\D{a}{y}A(y)$, with $a$ not free in $A$,
in case mode with the witness $0$: the premise
$\Gamma,A(0)\proves A(\Suc z)$ is the second hypothesis, weakened.
\rulename{DE} adds $\Q{y}A(y)$, and \rulename{QE} with the first hypothesis
gives $A(t)$.  If $A$ is an equation of $\mathfrak E$, pad the two
hypotheses by \cref{guard:lem:pad}, apply the compound case to $L(A(y))$,
and remove the padding of $L(A(t))$.
\end{proof}

The context of the two hypotheses is the same $\Gamma$ and does not contain
the case variable.  The lemma is not a rule for a case distinction on a
variable that occurs in an assumption.

\begin{lemma}[Implication between equations]
\label{guard:lem:imp}
For a compatible pair $(U,V)$ the following rules are admissible:
\[
 \infrule{\Gamma,\ U\proves V}{\Gamma\proves U\Rightarrow V},
 \qquad
 \infrule{\Gamma\proves U\Rightarrow V\qquad\Gamma\proves U}
         {\Gamma\proves V}.
\]
\end{lemma}

\begin{proof}
Write $\mu_w$ for $\mu_w(U,V)$.  The matrix $L(\mu_w)$ of $U\Rightarrow V$
is compound, so its introduction is in case mode, with the premise
$\Gamma,L(\mu_0)\proves L(\mu_{\Suc z})$.  Over $\Gamma,L(\mu_0)$:
\cref{guard:lem:pad} gives $\mu_0$, (G3) gives $U$, the given derivation and
cut give $V$, (G3) gives $\mu_{\Suc z}$, and \cref{guard:lem:pad} gives
$L(\mu_{\Suc z})$.  For the elimination, \rulename{DE} adds
$\Q{w}L(\mu_w)$; from $U$, (G3) and padding give $L(\mu_0)$; \rulename{QE}
at $1$ gives $L(\mu_1)$; and padding and (G3) give $V$.  The conversions
use transport outside the fresh binders only; no congruence of a binder is
assumed.
\end{proof}

\begin{theorem}[Equation induction from a guard]
\label{guard:thm:ei}
Let $E(y)$ be an equation of $\mathfrak E$ and $z$ a variable not free in
$\Gamma$ and $E$.  If $\Gamma\proves E(0)$ and
$\Gamma,E(z)\proves E(\Suc z)$, then $\Gamma\proves E(t)$ for every term
$t$.  The derivation uses the step premise at its own eigenvariable only:
no substitution into the given derivations is made.
\end{theorem}

\begin{proof}
Put
\[
 J(y)\defeq\gamma_yE(y),\qquad \Phi(y)\defeq\bigl(J(y)\Rightarrow E(y)\bigr),
\]
which is defined by (G4) and is a compound public formula.

\emph{The implication at every term.}  $\Gamma\proves\Phi(0)$ by the
introduction rule of \cref{guard:lem:imp}: the base $E(0)$ is derivable
over $\Gamma$, and the assumption $J(0)$ is not used.  For a fresh $w$,
$\Gamma\proves\Phi(\Suc w)$ by the same rule: the assumption $J(\Suc w)$ is
$\gamma_{\Suc w}E(\Suc w)$, which gives $E(\Suc w)$ by (G2).  Both
derivations are over the context $\Gamma$, which does not contain the case
variable, so \cref{guard:lem:nc} applies to the compound formula $\Phi$ and
gives $\Gamma\proves\Phi(v)$ for every term $v$.

\emph{Progress of the guard.}  Over $\Gamma,J(z)$ the elimination rule and
$\Phi(z)$ give $E(z)$; the step premise and cut give $E(\Suc z)$; and (G2)
gives $\gamma_{\Suc z}E(\Suc z)$, which is $J(\Suc z)$.  Thus
$\Gamma,J(z)\proves J(\Suc z)$.

\emph{Conclusion.}  By (G1) the matrix $J(y)$ is atomic and admissible, so
$\rulename{DI}_{\mathrm{root}}$ in progress mode, with an unused witness
coordinate $a$ and the witness $0$, gives $\Gamma\proves\D{a}{y}J(y)$.
\rulename{DE} adds $\Q{y}J(y)$.  The base $J(0)=\gamma_0E(0)$ is derivable by
(G2), and \rulename{QE} at $t$ gives $J(t)$.  With $\Phi(t)$ the elimination
rule gives $E(t)$.
\end{proof}

\begin{remark}[Two forms of the argument]
\label{guard:rem:ei-forms}
\pkg{guarded\_core}, \pkg{guarded\_orbit} and \pkg{protected\_core} use the
implication $J(y)\Rightarrow E(\Suc y)$ in place of $\Phi$.  Its two
branches then need the step premise at $0$ and at $\Suc w$, that is, a
substitution into the given derivation, and the conclusion needs a second
application of numerical cases.  That form is correct wherever substitution
is admissible: in $\mathsf G$ and $\mathsf H$ by \cref{guard:lem:struct},
and in $\mathsf G^{-}_{\mathrm{orb}}$ after the unrestricted transport has
been recovered, where \pkg{guarded\_orbit} performs the substitution in the
calculus with the derived rule as an inference and then expands that rule.
The form of the proof above is the one of \pkg{mirror\_guard}, where it is
needed because substitution is not available under the source restriction;
it applies to every scheme.  In each form the case distinction is made on a
compound implication over a context that does not contain the case variable;
a case distinction made while $J(z)$ is retained would violate the
eigenvariable condition.
\end{remark}

\subsection{The translation and its two composites}

\begin{definition}[Translation into a public language]
\label{guard:def:T}
Let a guard scheme on $\mathfrak E$ be given, and let $\varepsilon$ map
every equation to an equation of $\mathfrak E$, commuting with substitution.
The translation $T$ of native formulas is
\begin{align*}
 T(s=t) &\defeq \varepsilon(s=t),\\
 T(\Q{y}B) &\defeq \D{a}{y}P(TB), \qquad a \text{ fresh},\\
 T(\D{x}{y}B) &\defeq \D{x}{y}P(TB).
\end{align*}
\end{definition}

$T$ leaves every term unchanged outside $\varepsilon$, commutes with
substitution up to the names of bound variables, and maps every native
formula to a public formula: every translated matrix is padded, hence
compound, and the atomic matrices inside a padding are admissible by (G1).
The embedding $\varepsilon$ is the identity except for protected atoms,
where it is the wrapping $\varepsilon(s=t)=(W(s)=W(t))$.

\begin{theorem}[The public-language theorem]
\label{guard:thm:public}
Let $\mathsf K=\mathsf K(\mathcal F,\mathcal S)$,
$\mathcal P=\mathcal P(\mathfrak A,\mathfrak M)$ with
$\mathcal P^{Q}\subseteq\mathcal F$, a guard scheme on $\mathfrak E$ and an
embedding $\varepsilon$ be given, and assume:
\begin{enumerate}[label=\textup{(T\arabic*)}]
\item $\varepsilon(E)\dashv\vdash_{\Nat}E$ for every equation $E$;
\item $\mathsf K$ derives $\varepsilon(E)$ for every instance $E$ of
 \rulename{I0} and \rulename{IS}, and $\Gamma\proves\varepsilon(s=t)$ and
 $\Gamma\proves\varepsilon(h(s)=r)$ imply
 $\Gamma\proves\varepsilon(r=h(t))$;
\item $U\dashv\vdash_{\mathsf K}\varepsilon(U)$ for every equation
 $U\in\mathfrak A$;
\item $\Gamma\proves_{\mathsf K}B(t,0)$ for every matrix
 $B\in\mathfrak M$, every $\Gamma$ and every $t$.
\end{enumerate}
Then:
\begin{enumerate}[label=\textup{(\alph*)}]
\item $\Gamma\proves_{\Nat}A$ implies $T\Gamma\proves_{\mathsf K}TA$, for
 all native $\Gamma$ and $A$;
\item $A\dashv\vdash_{\Nat}TA$ for every native formula $A$;
\item $B\dashv\vdash_{\mathsf K}TB$ for every public formula $B$;
\item for public $\Gamma$ and $B$:
 $\Gamma\proves_{\mathsf K}B$ if and only if $\Gamma\proves_{\Nat}B$.
\end{enumerate}
The transformations of derivations are effective.
\end{theorem}

\begin{proof}
For a native matrix $B$, the translated closure is
$\mathrm{Cl}_yB\defeq\D{a}{y}P(TB)$, so $T(\Q{y}B)=\mathrm{Cl}_yB$.

\emph{(a).}  By induction on the derivation in $\Nat$.

\rulename{Hyp}: $TA\in T\Gamma$.

\rulename{I0}, \rulename{IS}, \rulename{Tr}: these are (T2); the rule
\rulename{Tr} of $\Nat$ infers $r=h(t)$ from $s=t$ and $h(s)=r$.

\rulename{QE}: the induction hypotheses are
$T\Gamma\proves\mathrm{Cl}_yB$ and $T\Gamma\proves TB(0)$.  \rulename{DE}
on the first adds the resource $\Q{y}P(TB)(y)$.  Padding gives $P(TB)(0)$,
\rulename{QE} at $t$ gives $P(TB)(t)$, and removal of the padding gives
$TB(t)=T(B(t))$.

\rulename{DE}: the induction hypotheses are
$T\Gamma\proves\D{x}{y}P(TB)$ and
$T\Gamma,\mathrm{Cl}_yB(b,y)\proves TA$.  \rulename{DE} in $\mathsf K$ on
the first adds $\Q{y}P(TB)(b,y)$.  Over this context derive
$\mathrm{Cl}_yB(b,y)$: its matrix is compound, the introduction is in case
mode with the witness $0$, and its premise
$P(TB)(b,0)\proves P(TB)(b,\Suc z)$ is \rulename{QE} on the resource and
the local base.  Cut $\mathrm{Cl}_yB(b,y)$ into the second hypothesis, which
gives $TA$, and close \rulename{DE}; $b$ is not free in $T\Gamma$ and $TA$.
This treats every use of the hypothesis $\Q{y}B(b,y)$ of the native
derivation, not only its uses as the major premise of \rulename{QE}.

\rulename{DI}: the induction hypothesis is
$T\Gamma,TB(t,0)\proves TB(t,\Suc z)$.  The matrix $P(TB)$ is compound, so
the introduction of $\D{x}{y}P(TB)$ is in case mode, with the premise
$T\Gamma,P(TB)(t,0)\proves P(TB)(t,\Suc z)$: remove the padding of the
base, cut it into the induction hypothesis, and pad the result.

\rulename{QI}: the induction hypothesis is
$T\Gamma,\varepsilon E(z)\proves\varepsilon E(\Suc z)$, and the goal is
$\D{a}{y}L(\varepsilon E(y))$.  Its introduction is in case mode, with the
premise $T\Gamma,L(\varepsilon E(0))\proves L(\varepsilon E(\Suc z'))$ for a
fresh $z'$.  Over $T\Gamma,L(\varepsilon E(0))$ the padding gives the base
$\varepsilon E(0)$; \cref{guard:thm:ei} for the equation
$\varepsilon E(y)\in\mathfrak E$, with the step premise weakened, gives
$\varepsilon E(\Suc z')$; and padding gives the goal.

\emph{(b).}  By induction on $A$, using three facts about $\Nat$:
$\Q{y}B\dashv\vdash\D{a}{y}B$ for $a$ not free in $B$
(\cref{base:lem:dummy}); interderivable matrices give interderivable
formulas under $\Q{}$ and under $\D{}{}$ (\cref{base:lem:bindercong}); and
$L(U)\dashv\vdash_{\Nat}U$.  The third holds because
(G2) holds in $\Nat$ by \cref{guard:lem:incl} and cut, and the proof of
\cref{guard:lem:pad} then applies in $\Nat$, where the introduction by cases
is available at every matrix.  The case of an equation is (T1).  For
$\Q{y}B$: $\Q{y}B\dashv\vdash\D{a}{y}B\dashv\vdash\D{a}{y}TB\dashv\vdash
\D{a}{y}P(TB)$.  The case of $\D{x}{y}B$ is the same without the first
step.

\emph{(c).}  By induction on the number of binders of $B$, for all
instances of $B$ at once.  For an equation of $\mathfrak A$ this is (T3).

Let $B=\D{x}{y}C$ with $C$ compound.  Then $TC$ is compound and
$TB=\D{x}{y}TC$.  From $\Gamma\proves\D{x}{y}C$, \rulename{DE} adds
$\Q{y}C(b,y)$; the introduction of $\D{x}{y}TC$ at the witness $b$ is in
case mode, and its premise $TC(b,0)\proves TC(b,\Suc z)$ is obtained by
converting the base to $C(b,0)$ with the induction hypothesis, applying
\rulename{QE} at $\Suc z$, and converting $C(b,\Suc z)$ to $TC(b,\Suc z)$.
The converse is symmetric.

Let $B=\D{x}{y}C$ with $C\in\mathfrak M$.  Then $TB=\D{x}{y}L(\varepsilon
C)$, and every instance of $C$ is an equation of $\mathfrak A$.  From
$\Gamma\proves\D{x}{y}C$, \rulename{DE} adds $\Q{y}C(b,y)$.  The
introduction of $\D{x}{y}L(\varepsilon C)$ at the witness $b$ is in case
mode; its local base is not used: $C(b,0)$ is derivable by (T4),
\rulename{QE} gives $C(b,\Suc z)$, (T3) gives $\varepsilon C(b,\Suc z)$,
and padding gives $L(\varepsilon C(b,\Suc z))$.  Conversely, from
$\Gamma\proves\D{x}{y}L(\varepsilon C)$, \rulename{DE} adds
$\Q{y}L(\varepsilon C(b,y))$.  The introduction of $\D{x}{y}C$ at the
witness $b$ is in progress mode; its local assumption $C(b,z)$ is not used:
$C(b,0)$ is derivable by (T4), (T3) and padding give
$L(\varepsilon C(b,0))$, \rulename{QE} gives $L(\varepsilon C(b,\Suc z))$,
and removal of the padding and (T3) give $C(b,\Suc z)$.  This direction uses
the provable base; it fails for an arbitrary atomic matrix, by
\cref{guard:prop:strict}.

\emph{(d).}  The implication from left to right is \cref{guard:lem:incl}.
Conversely let $\Gamma\proves_{\Nat}B$.  By (a),
$T\Gamma\proves_{\mathsf K}TB$.  By (c), $\Gamma\proves_{\mathsf K}TC$ for
every $C\in\Gamma$, so cut gives $\Gamma\proves_{\mathsf K}TB$, and (c)
gives $\Gamma\proves_{\mathsf K}B$.
\end{proof}

\begin{corollary}[Bi-interpretation]
\label{guard:cor:biinterp}
Under the hypotheses of \cref{guard:thm:public}, the theory of the public
sequents of $\mathsf K$, with the ordinary reading
$(\D{x}{y}B)^\circ=\exists x\,(B^\circ(x,0)\to\forall y\,B^\circ(x,y))$, is
bi-interpretable with $\iPRA$.  Both numerical domains are whole, equality
is equality, both numerical composites have the graph $x=y$, and both
formula composites are provably equivalent to the identity.  In particular
$\mathsf K$ derives its contradiction $\varepsilon(0=1)$ if and only if
$\iPRA$ proves $\bot$.
\end{corollary}

\begin{proof}
The inclusion of the public language into the native language and the
translation $T$ are the identity on the numerical sort and on equality; if
$\varepsilon$ is the wrapping, the identity graph is expressed by
$W(x)=W(y)$, which is $x=y$ in $\Nat$ by (T1).  They preserve derivability
by \cref{guard:lem:incl} and by (a), and their two composites are (b) and
(c).  Composition with the bi-interpretation of $\Nat$, $\Teq$ and $\iPRA$
(\cref{base:thm:bridge}) gives the statement.  For the contradiction:
$\varepsilon(0=1)$ is public, so by (d) it is derivable in $\mathsf K$ if
and only if it is derivable in $\Nat$, that is, by (T1), if and only if
$\Nat$ derives $0=1$.
\end{proof}

The arithmetic of the bridge, that is, the predecessor, the reflection of
equality, quantifier-free induction, pairing and the compiler for primitive
recursion (\cref{base:thm:pred,base:thm:reflect,base:thm:primrec}), is not
redone in any public language: it is transferred by (a) and (d).

\subsection{The guarded core}
\label{guard:sec:core}

This subsection draws on \pkg{guarded\_core}.  The guarded core is
$\mathsf G=\mathsf K(\mathcal P_0^{Q},\text{all terms})$: public formulas
are equations and paired formulas whose atomic matrices are zero-seeded,
transport is unrestricted, and the one introduction is determined by the
root.  Its seven rules are those of \cref{guard:def:calculus}.

\begin{theorem}[The guarded core]
\label{guard:thm:core}
\begin{enumerate}[label=\textup{(\roman*)}]
\item $\mathsf G$ admits weakening, substitution and cut.
\item The zero-seeded scheme on all equations is a guard scheme for
 $\mathsf G$ and $\mathcal P_0$.  Hence $\mathsf G$ admits numerical cases
 for every public formula, implication between equations, and contextual
 induction for every literal equation.
\item With $\varepsilon$ the identity, $T$ translates derivations of $\Nat$
 into derivations of $\mathsf G$; $A\dashv\vdash_{\Nat}TA$ for every native
 $A$; $B\dashv\vdash_{\mathsf G}TB$ for every $B\in\mathcal P_0$; and for
 $\Gamma,B$ in $\mathcal P_0$, $\Gamma\proves_{\mathsf G}B$ if and only if
 $\Gamma\proves_{\Nat}B$.
\item $\mathsf G$ on $\mathcal P_0$ is bi-interpretable and equiconsistent
 with $\iPRA$, with whole domains, equality as equality and identity
 numerical composites.
\end{enumerate}
\end{theorem}

\begin{proof}
(i) is \cref{guard:lem:struct}.  For (ii), all terms are sources, so
\cref{guard:lem:safe}(i),(ii) give the ordinary equality rules.  The two
sides of $G_0U$ are equal to $0$ and the two sides of $G_{\Suc w}(s=t)$ are
equal to $s$ and to $t$, by the selector equations; this is (G2).  The two
sides of $\mu_0(s=t,r=q)$ are equal to $s$ and $t$, and those of
$\mu_{\Suc w}$ to $r$ and $q$; this is (G3).  The consequences are
\cref{guard:lem:nc,guard:lem:imp,guard:thm:ei}.  For (iii) and (iv) the
hypotheses of \cref{guard:thm:public} hold: (T1) and (T3) are trivial; in
(T2) the computation equations are axioms and $r=h(t)$ follows from $s=t$
and $h(s)=r$ by transport in the template $h(v)=r$ and symmetry; (T4) is
\cref{guard:lem:zero}.  Then apply
\cref{guard:thm:public,guard:cor:biinterp}.
\end{proof}

\begin{corollary}[Administrative proofs are conservative]
\label{guard:cor:admin}
For $\Gamma$ and $B$ in $\mathcal P_0$,
\[
 \Gamma\proves_{\mathsf U}B
 \iff\Gamma\proves_{\mathsf G}B
 \iff\Gamma\proves_{\Nat}B .
\]
On $\mathcal P_0$ the mixed and the ordinary reading agree: $\Teq$ proves
$A^\dagger\leftrightarrow A^\circ$ for every $A\in\mathcal P_0$.
\end{corollary}

\begin{proof}
A derivation of $\mathsf G$ is a derivation of $\mathsf U$, a derivation of
$\mathsf U$ is transformed into one of $\Nat$ by \cref{guard:lem:incl}, and
\cref{guard:thm:core}(iii) closes the circle.  For the readings, at a
zero-seeded matrix $E(x,y)$ the instance $E(x,0)$ is provable, so with $x$
fixed $\forall y\,(E(x,y)\to E(x,\Suc y))$ and
$E(x,0)\to\forall y\,E(x,y)$ are equivalent: from the first, by induction
for the literal equation $E(x,y)$; from the second, because the base gives
every instance.  The equivalence passes under the existential quantifier,
and at a compound matrix the two readings have the same form, so induction
on $A$ applies.
\end{proof}

A derivation may therefore use unguarded atomic matrices and $\Q{}$
formulas as intermediate formulas without changing the public theory.  The
converse statement fails for $\mathsf U$ as a whole
(\cref{guard:prop:strict}): the full language of $\mathsf U$ is not claimed
to be bi-interpretable with $\iPRA$ under either reading.

\begin{remark}[Versions]
\label{guard:rem:core-versions}
\pkg{guarded\_core} states the comparison with the calculus of the reports
on the language of equations and $\D{}{}$ with clause contexts
(\cref{guard:rem:clauses}); the clause $T(\Q{y}B)=\D{a}{y}P(TB)$ is added in
\pkg{guarded\_orbit}, \pkg{deterministic\_auxiliary} and
\pkg{protected\_core}, and \pkg{guarded\_source} first replaces every
$\Q{y}B$ by $\D{a}{y}B$.  \pkg{guarded\_core} derives conservativity of
$\mathsf U$ from \cref{guard:lem:mixed-sound} and the agreement of the
readings, through $\iPRA$; the direct inclusion of \cref{guard:lem:incl} is
the argument of \pkg{guarded\_orbit} and \pkg{protected\_core}.  In Lean the
aim of the package, one root-determined introduction, is reached on the
full language by the witness-selected introduction
(\cref{kern:thm:witness}, \Lean{wprf\_iff\_nprf}); its own route is formalized in the common form
described in \cref{guard:rem:lean-public}.
\end{remark}

% ---------------------------------------------------------------------
\section{Which transport sources are needed}
\label{guard:sec:sources}

This section draws on \pkg{guarded\_source} for the dependency theorem; on
\pkg{deterministic\_auxiliary}, \pkg{guarded\_orbit} and \pkg{mirror\_guard},
each of which reproduces the same numerical argument inside its own
calculus; and on \pkg{guarded\_orbit} for the observation that the argument
needs induction at one equation only.  The singleton test, the filter and
the join at zero are taken by all four from \pkg{source\_dispatch}; they
are proved in \cref{kern:sec:dispatch} for the calculus $\mathsf Z$, which
has both introductions of $\Nat$, and are cited from there.  This section
proves what the guarded calculi need beyond that: the argument goes through
with two closure conditions in place of the introductions.

Transport with the sources $s\not\equiv 0$ is as strong as unrestricted
transport, on equations, as soon as a small amount of induction is
available.  The amount is made explicit, because in
\cref{guard:sec:composition} it is exactly what a strictly guarded language
fails to supply.

\begin{definition}[Restricted equational kernel]
\label{guard:def:rek}
A \emph{restricted equational kernel} is a relation $\Gamma\proves A$
between finite contexts and formulas of a language that contains all
equations, closed under weakening and under the rules \rulename{Hyp} for
equations, \rulename{I0}, \rulename{IS} and $\rulename{Tr}^{-}$ in equation
templates: from
$\Gamma\proves s=t$ with $s\not\equiv 0$ and $\Gamma\proves E(s)$ infer
$\Gamma\proves E(t)$, for a literal equation template $E(v)$.
\end{definition}

The derivability relation of every guarded calculus whose sources include
the terms $s\not\equiv 0$ is a restricted equational kernel, and
\cref{guard:lem:safe} holds for every restricted equational kernel, since
its proof uses these rules only.  Consider the following two closure
conditions on a restricted equational kernel.
\begin{enumerate}[label=\textup{(\Alph*)},start=3]
\item \emph{Cases at stable left sides.}  If $E(v)$ is an equation template
 whose left side is iterator-headed, $z$ is a variable not in $E$, and
 $\proves E(0)$ and $\proves E(\Suc z)$ in the empty context, then
 $\proves E(r)$ for every term $r$.
\end{enumerate}
\begin{enumerate}[label=\textup{(\Alph*)},start=15]
\item \emph{One step of an orbit.}  For every term $t\not\equiv 0$ and every
 context $\Gamma$, if $\Gamma\proves 0=\It{0}{t}{u}{\Suc u}$, then
 $\Gamma\proves 1=\It{1}{t}{u}{\Suc u}$.
\end{enumerate}

Define the terms
\begin{align*}
 \zeta(v)&\defeq\sel(v,1,0), & \kappa(v)&\defeq\sel(v,0,1), &
 \chi(v)&\defeq\sel(v,1,2),\\
 N(n)&\defeq\It{n}{0}{u}{\zeta(u)}, & R(n)&\defeq\It{n}{0}{u}{\chi(u)}, &
 \lambda(n)&\defeq N(R(n)).
\end{align*}
These are the terms $Z$, $K$, $H$, $N$, $R$ and $L$ of
\cref{kern:sec:dispatch}; three are renamed because $H_v$ and $L(\cdot)$
denote the mirror gate and the padding in this chapter.  The term $\zeta$
is the zero test, and $\lambda(\Suc v)$ tests whether $v$ is $0$.

\begin{lemma}[Tests and filter]
\label{guard:lem:tests}
In a restricted equational kernel with \textup{(C)}, the following equations
are derivable over every context, for every term $v$:
\begin{gather*}
 \chi(v)=\Suc\kappa(v),\qquad R(1)=1,\qquad R(\Suc\Suc v)=2,\qquad
 N(1)=1,\qquad N(2)=0,\\
 \lambda(1)=1,\qquad\lambda(\Suc\Suc v)=0,\qquad
 \lambda(\Suc v)=\zeta(v),\qquad
 \sel(\zeta(v),v,0)=v .
\end{gather*}
\end{lemma}

\begin{proof}
The first eight equations are \cref{kern:lem:ztest}.  The derivations given
there are made over the empty context and use \rulename{I0}, \rulename{IS},
the operations of \cref{guard:lem:safe}(i)--(iv) at sources that are
iterator-headed, successor-headed or the numerals $1$ and $2$, and
numerical cases at the two matrices $\chi(y)=\Suc\kappa(y)$ and
$\lambda(\Suc y)=\zeta(y)$.  Both matrices have an iterator-headed left
side, so the two case distinctions are instances of (C), and the
derivations are valid in every restricted equational kernel with (C).

The filter of \cref{kern:lem:zfilter} is the equation
$v=\sel(\zeta(v),v,0)$, proved by cases at a matrix whose left side is a
variable; (C) does not cover it.  In the orientation stated here the left
side is iterator-headed and (C) applies.  At $0$, replace $\zeta(0)$ by
$1$; then $\sel(1,0,0)=0$ by \rulename{IS}.  At $\Suc z$, replace
$\zeta(\Suc z)$ by $0$; then $\sel(0,\Suc z,0)=\Suc z$ by \rulename{I0}.
Both replacements are congruences at an iterator-headed source.
\end{proof}

\begin{theorem}[Equational source elimination]
\label{guard:thm:dependency}
Let $\proves$ be a restricted equational kernel with \textup{(C)} and
\textup{(O)}.  Then $\proves$ has unrestricted equality: over every context,
symmetry, transitivity and congruence hold for all equations, and
$\Gamma\proves s=t$ and $\Gamma\proves E(s)$ imply $\Gamma\proves E(t)$ for
every equation template $E(v)$ and every source $s$.  Both \textup{(C)} and
\textup{(O)} hold if $\proves$ is closed under contextual induction for
literal equations: from $\Gamma\proves E(0)$ and
$\Gamma,E(z)\proves E(\Suc z)$, with $z$ fresh, infer $\Gamma\proves E(r)$.
\end{theorem}

\begin{proof}
\emph{From induction to \textup{(C)} and \textup{(O)}.}  For (C), weaken the
derivation of $E(\Suc z)$ by the assumption $E(z)$ and apply induction; the
induction hypothesis is not used.  For (O), write
$A(n)\defeq\It{n}{t}{u}{\Suc u}$ and apply induction to the equation
$n=A(n)$ at the target $1$: the base is the hypothesis of (O), and the step
$\Gamma,z=A(z)\proves\Suc z=A(\Suc z)$ is the step of the proof of
\cref{kern:lem:zsym}, whose transports have the sources $z$, $A(\Suc z)$
and $\Suc A(z)$.

\emph{Symmetry at the source $0$.}  Let $\Gamma\proves 0=t$ with
$t\not\equiv 0$; the claim is $\Gamma\proves t=0$.  Read the proof of
\cref{kern:lem:zsym} with the given derivation of $0=t$ in place of the
hypothesis.  It derives $0=A(0)$ from $0=t$, uses its induction for the
single conclusion $1=A(1)$, which is (O), and derives $\zeta(t)=1$ from
$1=A(1)$ by \rulename{IS} and \cref{guard:lem:tests}.  With the filter in
the orientation of \cref{guard:lem:tests} the last step reads: put
$q\defeq\sel(\zeta(t),t,0)$; the filter gives $q=t$; the replacement of
$\zeta(t)$ by $1$ and \rulename{IS} give $q=\sel(1,t,0)=0$; reverse $q=t$
at the source $q$ and pass through $q$ to obtain $t=0$.

\emph{Unrestricted equality.}  Symmetry, congruence and transitivity at
every source follow from symmetry at the source $0$ and
\cref{guard:lem:safe}(ii),(iv) as in the proof of \cref{kern:thm:dispatch},
which uses nothing else; the cut made there is not needed, because symmetry
at the source $0$ has been obtained as a derived rule.  For an equation
template $a(v)=b(v)$, congruence gives $a(s)=a(t)$ and $b(s)=b(t)$, and
symmetry and transitivity turn $a(s)=b(s)$ into $a(t)=b(t)$.
\end{proof}

The proof uses induction in two places only: through (C), where the
induction hypothesis is ignored, and through (O), at the single equation
$n=\It{n}{t}{u}{\Suc u}$ with $t$ and the context fixed.  The variable
source $z$ of the step of (O) is never replaced by $0$: the target $1$ is
supplied by the induction rule, not by a substitution into the step.  No
logical rule occurs: no binder, no introduction by cases.  This removes the
apparent circularity between the two reductions, in which the recovery of
transport at the source $0$ would need logical cases and the recovery of
logical cases at an equation would need unrestricted transport.

\begin{corollary}[The equational fragment]
\label{guard:cor:eminus}
Let $\mathcal E^{-}$ be the calculus whose formulas are the equations and
whose rules are \rulename{Hyp}, \rulename{I0}, \rulename{IS},
$\rulename{Tr}^{-}$ in equation templates and contextual induction for
literal equations.  Unrestricted transport in equation templates is
admissible in $\mathcal E^{-}$.  The restriction $s\not\equiv 0$ changes
neither its theorems nor its consequences from hypotheses.
\end{corollary}

\Cref{kern:thm:dispatch} is the case of \cref{guard:thm:dependency} for the
calculus $\mathsf Z$, which has \rulename{QI} and hence induction; what
\cref{guard:thm:dependency} adds is that the two closure conditions
suffice, which is what the calculi of \cref{guard:sec:composition} need.
The Lean statements for the fragment, \Lean{ZKernel.trF} and
\Lean{zprf\_iff\_nprf}, are recorded in \cref{kern:rem:dispatch-sources}.

\begin{proposition}[Transport in public templates]
\label{guard:prop:pubtr}
Let $\mathsf K$ be a guarded calculus whose proof language contains
$\mathcal P^{Q}$ for a public language
$\mathcal P=\mathcal P(\mathfrak A,\mathfrak M)$.  Assume that $\mathsf K$
has unrestricted equality in the sense of \cref{guard:thm:dependency}, and
that for every $B\in\mathfrak M$ either $\Gamma\proves B(t,0)$ for all
$\Gamma$ and $t$, or $\Gamma\proves\D{x}{y}B'$ for all $\Gamma$ and all
substitution instances $B'$ of $B$.  Then for every template
$F(v)\in\mathcal P$,
\[
 \Gamma\proves s=t\quad\text{and}\quad\Gamma\proves F(s)
 \quad\text{imply}\quad\Gamma\proves F(t).
\]
\end{proposition}

\begin{proof}
By induction on the number of binders of $F$, for all contexts and all
$s,t$.  For an equation this is the hypothesis.  Let
$F(v)=\D{x}{y}B(v,x,y)$.  If every instance of $\D{x}{y}B$ is derivable,
nothing is to be done.  Otherwise apply \rulename{DE} to $F(s)$, which adds
$\Q{y}B(s,b,y)$, and introduce $F(t)$ at the witness $b$.

If $B$ is compound the introduction is in case mode.  From its base
$B(t,b,0)$ the induction hypothesis, with the reversed equation $t=s$, gives
$B(s,b,0)$; \rulename{QE} gives $B(s,b,\Suc z)$; and the induction
hypothesis gives $B(t,b,\Suc z)$.

If $B$ is atomic the introduction is in progress mode and its local
assumption is not used: $B(s,b,0)$ is derivable, \rulename{QE} gives
$B(s,b,\Suc z)$, and transport in the equation template gives
$B(t,b,\Suc z)$.
\end{proof}

Hence, under the hypotheses of the proposition, a calculus
$\mathsf K(\mathcal P^{Q},\,s\not\equiv 0)$ derives the same sequents as
$\mathsf K(\mathcal P^{Q},\,\text{all terms})$: every unrestricted transport
is replaced by its derivation and the other inferences are kept.  The
proposition is about public templates.  No transport at the source $0$ is
asserted in a template that contains an unguarded atomic matrix below a
binder.

\begin{remark}[The crossed rule]
\label{guard:rem:crossed}
The induction of \cref{guard:thm:dependency} is the price of the uncrossed
conclusion.  For the rule \rulename{Tr} of the kernel the restriction
$s\not\equiv 0$ costs nothing: \cref{kern:lem:oseq}(ii) reverses an
equation with the source $0$ in one inference, by a constant template, and
needs neither the tests nor an induction; the same holds for the reversing
rule of \cref{kern:thm:reversal}.  The calculi of
\crefrange{guard:sec:shapes}{guard:sec:composition} have the uncrossed
rule, and for it the obstruction of \cref{guard:prop:inert} is real.
\end{remark}

\begin{remark}[Substitution under a source restriction]
\label{guard:rem:substwording}
As for the calculi of \cref{kern:rem:ossubst}, substitution is admissible
in a source-restricted guarded calculus once its equivalence with the
unrestricted calculus is known, and what fails is the rule-by-rule
substitution of \cref{guard:lem:struct}; \pkg{guarded\_source},
\pkg{guarded\_orbit} and \pkg{protected\_core} make the same point.  The
equivalence must be proved first, and none of the proofs of this chapter
substitutes into a derivation with $\rulename{Tr}^{-}$ before that point.
\end{remark}

% ---------------------------------------------------------------------
\section{The failure of composition and its repairs}
\label{guard:sec:composition}

This section draws on \pkg{mirror\_guard}, \pkg{guarded\_orbit} and
\pkg{deterministic\_auxiliary} for the obstruction; on \pkg{guarded\_source}
and \pkg{deterministic\_auxiliary} for the repair by auxiliary formulas; on
\pkg{guarded\_orbit} for the repair by the orbit; on \pkg{mirror\_guard} for
the mirror guard; and on \pkg{protected\_core} for protected atoms.

The guarded core removes the choice of the introduction mode
(\cref{guard:thm:core}), and the source restriction removes the transports
at the literal $0$ (\cref{guard:thm:dependency}).  The first result is
stated for unrestricted transport and the second for a calculus with
induction at arbitrary equations.  In the strictly guarded calculus
$\mathsf G^{-}=\mathsf K(\mathcal P_0^{Q},\,s\not\equiv 0)$ neither
hypothesis is available, and the two results do not compose.

\subsection{The obstruction}

An equation is \emph{zero-left} if its left side is the literal $0$.  A
formula of a proof language $\mathcal P^{Q}$ is \emph{safe} if it is not a
zero-left equation: it is a public compound formula, a resource, or an
equation whose left side is not the literal $0$.

\begin{proposition}[An inert zero-left hypothesis]
\label{guard:prop:inert}
Let $\mathcal P=\mathcal P(\text{all equations},\mathfrak M)$, where the
left side of every matrix of $\mathfrak M$ is iterator-headed or
successor-headed, and let
$\mathsf K=\mathsf K(\mathcal P^{Q},\,s\not\equiv 0)$.  If $\Gamma$ consists
of safe formulas, $F$ is safe and $\Gamma,\,0=t\proves_{\mathsf K}F$, then
$\Gamma\proves_{\mathsf K}F$.
\end{proposition}

\begin{proof}
By induction on the derivation, for all safe $\Gamma$ and $F$.  Every
instance of a matrix of $\mathcal P$ is safe: a compound formula stays
compound, and the left side of an instance of a matrix of $\mathfrak M$
keeps its head.

\rulename{Hyp}: the conclusion is safe, so it is not the hypothesis $0=t$
and belongs to $\Gamma$.  \rulename{I0}, \rulename{IS}: no hypothesis is
used.

$\rulename{Tr}^{-}$, with the premises $p=q$ and $F(p)$ and the safe
conclusion $F(q)$: the major premise is safe because $p\not\equiv 0$.  The
minor premise is safe: if $F$ is compound this is clear; if $F(v)$ is
$a(v)=b(v)$ and $a(p)\equiv 0$, then $a(v)$ is the constant $0$, because a
constructor-headed term and the variable $v$ do not become the literal $0$
when $v$ is replaced by $p\not\equiv 0$; then $a(q)\equiv 0$, against the
safety of $F(q)$.  The induction hypothesis applies to both premises.

\rulename{QE}, with the premises $\Q{y}B(y)$ and $B(0)$: both are safe, the
first as a resource and the second as an instance of a matrix.

\rulename{DE}: the major premise is compound; the minor premise has the
same conclusion and one more hypothesis, the resource $\Q{y}B(b,y)$, which
is safe.

$\rulename{DI}_{\mathrm{root}}$: the local assumption $B(t',\eta_B(z))$ and
the conclusion $B(t',\Suc z)$ of the premise are instances of a matrix,
hence safe.
\end{proof}

\begin{corollary}[The strictly guarded composition is incomplete]
\label{guard:cor:obstruction}
In $\mathsf G^{-}$ the sequents $0=1\proves 1=0$ and $0=x\proves x=0$, for a
variable $x$, are not derivable.  Both are derivable in $\mathsf G$.  The
same holds for every calculus of \cref{guard:prop:inert}, for instance with
the atomic matrices $\sel(y,t,s)=t$.
\end{corollary}

\begin{proof}
The conclusions $1=0$ and $x=0$ are safe.  By \cref{guard:prop:inert} a
derivation would give $\proves 1=0$, respectively $\proves x=0$, in the
empty context, hence in $\Nat$ by \cref{guard:lem:incl}; both fail in the
natural numbers.  In $\mathsf G$ symmetry is \cref{guard:lem:safe}(ii).
\end{proof}

The proposition says where the composition breaks.  Every equation that a
strictly guarded derivation can obtain from a binder has a stable left
side: the instances of a zero-seeded matrix, the local assumptions and the
premises of the introductions.  A transport with a source $s\not\equiv 0$
cannot turn a zero-left equation into a safe one.  A zero-left hypothesis
therefore has no path to a safe conclusion.  The proposition does not apply
to a context with a clause whose antecedent is zero-left and whose
consequent is safe, and each of the repairs below supplies such a clause or
removes the zero-left hypotheses from the public language.

\begin{remark}[Versions of the obstruction]
\label{guard:rem:obstruction-versions}
\Cref{guard:prop:inert} is stated in \pkg{mirror\_guard}, for every guard
with a constructor-headed left side, and in \pkg{guarded\_orbit} for the
hypothesis $0=x$.  \pkg{deterministic\_auxiliary} proves the first sequent
of \cref{guard:cor:obstruction} underivable by a counterinterpretation over
the natural numbers: every zero-left equation is read as true, every other
equation as equality, and the compound formulas by their ordinary reading.
Every equation below a binder of a formula of $\mathcal P_0^{Q}$ has an
iterator-headed left side, so it is read as equality and the reading
commutes with substitution; the computation equations, both modes of the
introduction and \rulename{QE} and \rulename{DE} are sound; and
$\rulename{Tr}^{-}$ is sound because its major premise $p=q$ is not
zero-left and an equation template whose instance at $p$ is zero-left has
the constant left side $0$, so that its instance at $q$ is true by
convention.  The hypothesis $0=1$ is true and the conclusion $1=0$ is false.
In Lean the calculus \Lean{PGPrf} with the root-determined introduction and
transport at iterator-headed sources does not derive $x=z$ from $x=y$ and
$y=z$ between variables (\Lean{pgprf\_not\_trans}), by a reading in which a
value is a number with a colour and equality is not transitive; hence it is
strictly weaker than the kernel on arbitrary sequents
(\Lean{pgprf\_strictly\_weaker}).
\end{remark}

\subsection{First repair: auxiliary formulas}
\label{guard:sec:aux}

This subsection draws on \pkg{guarded\_source} and
\pkg{deterministic\_auxiliary}.  The public language remains
$\mathcal P_0$, and derivations are made in
$\mathsf U^{-}=\mathsf K(\mathcal L,\,s\not\equiv 0)$: the proof language is
the full native language, in which an atomic matrix below $\D{}{}$ need not
be guarded.  A public theorem is a public formula with a derivation in
$\mathsf U^{-}$; its intermediate formulas need not be public.  This
convention is part of the calculus.  By \cref{guard:rem:se} its rules are
\rulename{Hyp}, the two computation equations, \rulename{SE}, \rulename{DE}
and $\rulename{DI}_{\mathrm{root}}$.

\begin{theorem}[Source selection and root-determined introduction, with
auxiliary formulas]
\label{guard:thm:aux}
\begin{enumerate}[label=\textup{(\roman*)}]
\item $\mathsf U^{-}$ admits contextual induction for every literal
 equation, by a derivation without transport, and numerical cases for every
 native formula.
\item $\mathsf U^{-}$ has unrestricted equality and admits transport, at
 every source, in every template of $\mathcal P_0$.
\item With the zero-seeded scheme and $\varepsilon$ the identity, $T$
 translates derivations of $\Nat$ into derivations of $\mathsf U^{-}$, and
 $B\dashv\vdash_{\mathsf U^{-}}TB$ for every $B\in\mathcal P_0$.
\item For $\Gamma$ and $B$ in $\mathcal P_0$,
 \[
  \Gamma\proves_{\mathsf U^{-}}B
  \iff\Gamma\proves_{\mathsf G}B
  \iff\Gamma\proves_{\Nat}B ,
 \]
 and $\mathsf U^{-}$ on $\mathcal P_0$ is bi-interpretable and
 equiconsistent with $\iPRA$, with whole domains, equality as equality and
 identity numerical composites.
\end{enumerate}
\end{theorem}

\begin{proof}
(i) Let $\Gamma\proves E(0)$ and $\Gamma,E(z)\proves E(\Suc z)$.  The
formula $\D{a}{y}E(y)$, with $a$ not free in $E$, belongs to $\mathcal L$,
and its introduction is in progress mode: with the witness $0$ its premise
is the given step.  \rulename{DE} adds $\Q{y}E(y)$, in which the
eigenvariable does not occur, and \rulename{QE} with the base gives $E(t)$.
For numerical cases at an equation, use this induction and ignore the
induction hypothesis; at a compound formula $A(y)$, introduce
$\D{a}{y}A(y)$ in case mode and proceed in the same way.

(ii) By (i) and \cref{guard:thm:dependency}, $\mathsf U^{-}$ has
unrestricted equality; additional hypotheses of the context are passive in
that proof.  \Cref{guard:prop:pubtr} with \cref{guard:lem:zero} gives the
transport in public templates.

(iii) and (iv).  With unrestricted equality the zero-seeded scheme
satisfies (G2) and (G3) in $\mathsf U^{-}$ exactly as in
\cref{guard:thm:core}, and (T1) to (T4) hold for the same reasons.  Apply
\cref{guard:thm:public,guard:cor:biinterp}; the middle member of the chain
is \cref{guard:thm:core}(iii).
\end{proof}

The induction of (i) is made at an arbitrary literal equation, and the
matrix $E$ is in general not guarded.  The proof of
\cref{guard:thm:dependency} uses it at the equation
$n=\It{n}{t}{u}{\Suc u}$, whose left side is a variable.  By
\cref{guard:cor:obstruction} this use of an unguarded matrix cannot be
replaced by guarded ones: the auxiliary formulas are a necessary part of the
calculus, and they are why the proof language is stated explicitly.  The
theorem asserts the equivalence and the composites for the public language.
It asserts neither a formula composite nor transport at the source $0$ for
all auxiliary formulas; by \cref{guard:prop:strict} the former fails.

\begin{remark}[Versions]
\label{guard:rem:aux-versions}
\pkg{guarded\_source} proves \cref{guard:thm:dependency} for the fragment
$\mathcal E^{-}$ and instantiates it in $\mathsf U^{-}$, with the filter in
the orientation $v=\sel(\zeta(v),v,0)$; \pkg{deterministic\_auxiliary} runs
the numerical argument directly in $\mathsf U^{-}$, with the filter as in
\cref{guard:lem:tests}.  Both prove the implication from $\mathsf U^{-}$ to
$\mathsf G$ in (iv) through $\iPRA$, by \cref{guard:lem:mixed-sound}, the
agreement of the readings on $\mathcal P_0$ (\cref{guard:cor:admin}) and the
faithfulness of the bridge; \cref{guard:lem:incl} gives it directly.  In
Lean the combination of source selection with a single introduction is
proved on the full native language, with the witness-selected introduction
and the crossed rule (the seven-rule calculus of \cref{kern:def:OM},
\Lean{mprf\_iff\_nprf}); the auxiliary formulas of the
two packages are not formalized.
\end{remark}

\subsection{Second repair: the administrative orbit}
\label{guard:sec:orbit}

This subsection draws on \pkg{guarded\_orbit}.  The auxiliary formulas of
the first repair are reduced to one family.  The proof language is
$(\mathcal P_0^{+})^{Q}$: the guarded language together with the paired
formulas over orbit matrices $y=\It{y}{a}{u}{\Suc u}$, $y$ not free in $a$.
The calculus is
$\mathsf G^{-}_{\mathrm{orb}}=\mathsf K((\mathcal P_0^{+})^{Q},\,s\not\equiv
0)$, with the rules of \cref{guard:def:calculus}; no rule is added for the
orbit matrices.  The public language remains $\mathcal P_0$.

An equation is \emph{stable-left} if its left side is iterator-headed or
successor-headed.  The class of stable-left equations is closed under
substitution, and none of its members is zero-left.

\begin{theorem}[The orbit bootstrap]
\label{guard:thm:orbit}
\begin{enumerate}[label=\textup{(\roman*)}]
\item In $\mathsf G^{-}$, hence in $\mathsf G^{-}_{\mathrm{orb}}$, the
 zero-seeded scheme on the stable-left equations is a guard scheme for
 $\mathcal P_0$.  Hence numerical cases hold for every compound public
 formula and every stable-left equation, and condition \textup{(C)} of
 \cref{guard:sec:sources} holds.
\item For every term $t$ with $y$ not free in it and every context,
 $\mathsf G^{-}_{\mathrm{orb}}$ derives the \emph{orbit closure}
 \[
  \D{a}{y}\bigl(y=\It{y}{t}{u}{\Suc u}\bigr),
 \]
 and condition \textup{(O)} holds.
\item $\mathsf G^{-}_{\mathrm{orb}}$ has unrestricted equality, admits
 transport at every source in every template of $\mathcal P_0^{+}$, and
 derives the same sequents as
 $\mathsf K((\mathcal P_0^{+})^{Q},\text{all terms})$.
\item With the zero-seeded scheme on all equations and $\varepsilon$ the
 identity, $T$ translates derivations of $\Nat$ into derivations of
 $\mathsf G^{-}_{\mathrm{orb}}$; $B\dashv\vdash TB$ for every
 $B\in\mathcal P_0$; and for $\Gamma$ and $B$ in $\mathcal P_0$,
 \[
  \Gamma\proves_{\mathsf G^{-}_{\mathrm{orb}}}B
  \iff\Gamma\proves_{\mathsf G}B
  \iff\Gamma\proves_{\Nat}B .
 \]
 The public theory is bi-interpretable and equiconsistent with $\iPRA$,
 with whole domains, equality as equality and identity numerical
 composites.
\end{enumerate}
\end{theorem}

\begin{proof}
(i) (G1) and (G4) hold by inspection, and $G_vU$ and $\mu_w(U,V)$ are
stable-left because their left sides are iterator-headed.  For (G2): the two
sides of $G_0U$ are iterator-headed and equal to $0$ by \rulename{I0}, so
\cref{guard:lem:safe}(iv) proves $G_0U$.  For $U\equiv(s=t)$ put
$a\defeq\sel(\Suc w,0,s)$ and $b\defeq\sel(\Suc w,0,t)$; then $a=s$ and
$b=t$ by \rulename{IS}, the term $s$ is not the literal $0$ because $U$ is
stable-left, and \cref{guard:lem:safe}(v) gives $U\dashv\vdash G_{\Suc w}U$.
(G3) is the same lemma: for $\mu_0(s=t,r=q)$ with $a\defeq\sel(0,s,r)$,
$b\defeq\sel(0,t,q)$, and for $\mu_{\Suc w}$ with
$a\defeq\sel(\Suc w,s,r)$, $b\defeq\sel(\Suc w,t,q)$, where $r\not\equiv 0$
because the second equation is stable-left.  The consequences are
\cref{guard:lem:nc}, and (C) is its case of an equation with an
iterator-headed left side over the empty context.

(ii) Write $A(n)\defeq\It{n}{t}{u}{\Suc u}$.  The matrix $y=A(y)$ is an
orbit matrix, so the formula belongs to $\mathcal P_0^{+}$ and its
introduction is in progress mode.  With the witness $0$ the premise is
$\Gamma,z=A(z)\proves\Suc z=A(\Suc z)$, which is the step derived in the
proof of \cref{guard:thm:dependency}: the source $z$ is a variable.  No base
is needed.  For (O), let $\Gamma\proves 0=A(0)$.  \rulename{DE} on the orbit
closure adds $\Q{y}(y=A(y))$, and \rulename{QE} with the base at the target
$1$ gives $1=A(1)$.

(iii) By (i), (ii) and \cref{guard:thm:dependency} the calculus has
unrestricted equality.  \Cref{guard:prop:pubtr} applies: a zero-seeded
matrix has provable instances at zero (\cref{guard:lem:zero}), and every
paired formula over an orbit matrix is derivable by (ii).

(iv) With unrestricted equality the zero-seeded scheme on all equations
satisfies (G2) and (G3) as in \cref{guard:thm:core}, and (T1) to (T4) hold.
Apply \cref{guard:thm:public,guard:cor:biinterp} to the calculus
$\mathsf G^{-}_{\mathrm{orb}}$ and the public language $\mathcal P_0$.
\end{proof}

In the derivation of $0=t\proves t=0$ the orbit closure is eliminated with
the zero-left base $0=A(0)$ and yields the stable-left equation $1=A(1)$.
This is the clause that \cref{guard:prop:inert} excludes from a strictly
guarded context, and it is the only use of an unguarded matrix: by (i) the
cases needed by \cref{guard:lem:tests} are available in $\mathsf G^{-}$
itself.  The filter must be taken in the orientation of
\cref{guard:lem:tests}, with the iterator on the left; the equation
$v=\sel(\zeta(v),v,0)$ is not stable-left.

\begin{remark}[Scope and versions]
\label{guard:rem:orbit-versions}
The theorem also holds with $\mathcal P_0^{+}$ as the public language: for
an orbit matrix $O$ both $\D{x}{y}O$ and its translation $\D{x}{y}L(O)$ are
derivable, the second in case mode from its base by the elimination of the
orbit closure, so the two are interderivable.  The mixed reading of
\cref{guard:def:mixed} agrees with the ordinary reading on
$\mathcal P_0^{+}$ as well: at an orbit matrix both readings are provable,
the progressive one by the step of (ii) and the conditional one from it by
induction, so their agreement does not depend on a provable base.
\pkg{guarded\_orbit} states
the calculus with the full native language as proof language, that is
$\mathsf U^{-}$, and observes that the derivations constructed lie in
$\mathsf G^{-}_{\mathrm{orb}}$; it proves the induction of
\cref{guard:thm:ei} in the form discussed in \cref{guard:rem:ei-forms}.
The named derivation of $0=t\proves t=0$ recorded by the package has
$326$ distinct nodes, of which $171$ are transports, $122$ computation
equations and $10$ introductions, exactly one of them at an orbit matrix.
\end{remark}

\subsection{Third repair: the mirror guard}
\label{guard:sec:mirror}

This subsection draws on \pkg{mirror\_guard}.  The guard is turned around:
an atomic matrix below $\D{}{}$ has the mirrored shape $a=\It{y}{a}{u}{h}$,
with one iterator, the seed repeated on the left.  The calculus is
$\mathsf M=\mathsf K(\mathcal P_m^{Q},\,s\not\equiv 0)$ and is strictly
guarded: every formula of a derivation is public or a resource.  The left
side of a mirrored matrix is arbitrary, so the hypothesis of
\cref{guard:prop:inert} fails, and a zero-left equation $0=t$ has the mirror
gate $0=\sel(v,0,t)$, which is again zero-left: it can be padded without a
transport at the source $0$.

\begin{theorem}[The mirror guard]
\label{guard:thm:mirror}
\begin{enumerate}[label=\textup{(\roman*)}]
\item The mirror scheme on all equations is a guard scheme for $\mathsf M$
 and $\mathcal P_m$.  Hence $\mathsf M$ admits numerical cases for every
 public formula, implication between equations with a common left side, and
 contextual induction for every literal equation, by derivations that
 substitute into no given derivation.
\item $\mathsf M$ has unrestricted equality, admits transport at every
 source in every template of $\mathcal P_m$, and derives the same sequents
 as $\mathsf K(\mathcal P_m^{Q},\text{all terms})$.
\item With the mirror scheme and $\varepsilon$ the identity, $T$ translates
 derivations of $\Nat$ into derivations of $\mathsf M$;
 $A\dashv\vdash_{\Nat}TA$ for every native $A$;
 $B\dashv\vdash_{\mathsf M}TB$ for every $B\in\mathcal P_m$; and for
 $\Gamma,B$ in $\mathcal P_m$, $\Gamma\proves_{\mathsf M}B$ if and only if
 $\Gamma\proves_{\Nat}B$.
\item $\mathsf M$ on $\mathcal P_m$ is bi-interpretable and equiconsistent
 with $\iPRA$, with whole domains, equality as equality and identity
 numerical composites.
\end{enumerate}
\end{theorem}

\begin{proof}
(i) (G1) and (G4) hold by inspection.  (G2): $H_0(s=t)$ is
$s=\sel(0,s,t)$, the reverse of an instance of \rulename{I0}
(\cref{guard:lem:safe}(iii)).  For $H_{\Suc w}(s=t)$ put
$q\defeq\sel(\Suc w,s,t)$; then $q\not\equiv 0$ and $q=t$ by
\rulename{IS}, and \cref{guard:lem:safe}(vi) gives
$(s=q)\dashv\vdash(s=t)$.  (G3): for $\mu_0(s=a,s=b)$ take
$q\defeq\sel(0,a,b)$ with $q=a$, and for $\mu_{\Suc w}$ take
$q\defeq\sel(\Suc w,a,b)$ with $q=b$; \cref{guard:lem:safe}(vi) gives
$(s=q)\dashv\vdash(s=a)$ and $(s=q)\dashv\vdash(s=b)$.  Only the right side
of an equation is expanded or contracted; the left side $s$ is never
touched, and it may be the literal $0$ or a variable.  The consequences are
\cref{guard:lem:nc,guard:lem:imp,guard:thm:ei}.

(ii) By (i), $\mathsf M$ is closed under induction for literal equations,
so \cref{guard:thm:dependency} gives unrestricted equality.  A mirrored
matrix has provable instances at zero (\cref{guard:lem:zero}), so
\cref{guard:prop:pubtr} gives the transport in public templates and the
equality of the two consequence relations.

(iii) and (iv).  (G2) and (G3) hold by (i); (T1) and (T3) are trivial; (T2)
holds by unrestricted equality; (T4) is \cref{guard:lem:zero}.  Apply
\cref{guard:thm:public,guard:cor:biinterp}.
\end{proof}

In the induction of \cref{guard:thm:dependency} for the equation
$n=\It{n}{t}{u}{\Suc u}$, the guard $J(n)$ of \cref{guard:thm:ei} is
$n=\sel\bigl(n,n,\It{n}{t}{u}{\Suc u}\bigr)$: a mirrored matrix whose left
side is the induction variable.  The zero-seeded guard has no such matrix,
which is the content of \cref{guard:prop:inert}; the orbit matrix of the
second repair supplies one by hand.  The order of the recovery is: equality
at admissible sources; expansion on the right; padding and cases;
implication with a common left side; induction; equality at the source $0$;
transport in public templates; the translation.  The induction is recovered
before the source restriction is lifted, and it uses the implication
$J(y)\Rightarrow E(y)$, not $J(y)\Rightarrow E(\Suc y)$
(\cref{guard:rem:ei-forms}).

\begin{remark}[Scope]
\label{guard:rem:mirror-scope}
The mirrored shape shows one iterator where the zero-seeded shape shows
two, and it repeats the seed; no statement about the length of expanded
terms is made.  \pkg{mirror\_guard} states the calculus with clause contexts
and proves its theorem against the calculus of the reports on the language
of equations and $\D{}{}$ (\cref{guard:rem:clauses,guard:rem:core-versions}).
\end{remark}

\subsection{Fourth repair: protected atoms}
\label{guard:sec:protected}

This subsection draws on \pkg{protected\_core}.  The zero-left equations are
removed from the public language.  A public atom is a protected equation
$p=q$, both sides iterator-headed, and ordinary equality is expressed by
the wrapped equation:
\[
 s\simeq t\ \defeq\ \bigl(W(s)=W(t)\bigr),\qquad
 \bot_{W}\ \defeq\ \bigl(W(0)=W(1)\bigr).
\]
The public language is $\mathcal P_W$, with zero-seeded atomic matrices,
and the calculus is
$\mathsf H=\mathsf K(\mathcal P_W^{Q},\text{iterator-headed terms})$: its
seven rules are \rulename{Hyp}, \rulename{I0}, \rulename{IS},
$\rulename{Tr}_I$, \rulename{QE}, \rulename{DE} and
$\rulename{DI}_{\mathrm{root}}$.  Its proof language contains all
equations as auxiliary formulas, which are needed for the computation
equation $W(t)=t$ when $t$ is not iterator-headed; they are not public
atoms.  The source condition is stronger than $s\not\equiv 0$: a variable is
not a source.  In return the class of sources is closed under substitution.

\begin{lemma}[Lifting a wrapped equation]
\label{guard:lem:wlift}
For all terms $s,t$ and every term template $f(v)$, in $\mathsf H$,
\[
 \Gamma\proves W(s)=W(t)\quad\text{implies}\quad
 \Gamma\proves W(f(s))=W(f(t)).
\]
\end{lemma}

\begin{proof}
Congruence along $W(s)=W(t)$ in the template $W(f(v))$ gives
$W(f(W(s)))=W(f(W(t)))$.  Congruence along the computation equations
$W(s)=s$ and $W(t)=t$ in the same template gives
$W(f(W(s)))=W(f(s))$ and $W(f(W(t)))=W(f(t))$.  All sources and all six
sides are iterator-headed, so \cref{guard:lem:safe}(ii) reverses the first
of the last two equations and composes the three.  The argument does not
depend on whether $s$ is the literal $0$.
\end{proof}

\begin{theorem}[The protected core]
\label{guard:thm:protected}
\begin{enumerate}[label=\textup{(\roman*)}]
\item $\mathsf H$ admits weakening, substitution and cut.
\item The zero-seeded scheme on the protected equations is a guard scheme
 for $\mathsf H$ and $\mathcal P_W$.  Hence $\mathsf H$ admits numerical
 cases for every public formula, implication between protected equations,
 and contextual induction for every protected equation.
\item With $\varepsilon(s=t)\defeq(W(s)=W(t))$, $T$ translates derivations
 of $\Nat$ into derivations of $\mathsf H$, all of whose transports have
 iterator-headed sources; $A\dashv\vdash_{\Nat}TA$ for every native $A$;
 $B\dashv\vdash_{\mathsf H}TB$ for every $B\in\mathcal P_W$; and for
 $\Gamma,B$ in $\mathcal P_W$, $\Gamma\proves_{\mathsf H}B$ if and only if
 $\Gamma\proves_{\Nat}B$.
\item $\mathsf H$ on $\mathcal P_W$ is bi-interpretable and equiconsistent
 with $\iPRA$.  Both numerical domains are whole; equality is equality,
 expressed in the public language by $\simeq$; both numerical composites
 are the identity, with the graph $x\simeq y$; and
 \[
  \proves_{\mathsf H}\bot_W\iff\proves_{\Nat}0=1\iff\iPRA\proves\bot .
 \]
\end{enumerate}
\end{theorem}

\begin{proof}
(i) is \cref{guard:lem:struct}: an iterator-headed source keeps its head
under substitution, and the zero-seeded shape is preserved.  The argument
would fail for the sources $s\not\equiv 0$.

(ii) (G1) and (G4) hold by inspection, and the gates and selectors of
protected equations are protected.  (G2): the two sides of $G_0U$ are
iterator-headed and equal to $0$, so \cref{guard:lem:safe}(iv) proves
$G_0U$; and for $U\equiv(p=q)$, \cref{guard:lem:safe}(v) with
$a\defeq\sel(\Suc w,0,p)$ and $b\defeq\sel(\Suc w,0,q)$ gives
$U\dashv\vdash G_{\Suc w}U$, since $p$ and $q$ are sources.  (G3) is the
same lemma applied to $\sel(0,p,r)$, $\sel(0,q,s)$ and to
$\sel(\Suc w,p,r)$, $\sel(\Suc w,q,s)$.

(iii) and (iv).  (T1): $W(s)=s$ and $W(t)=t$ in $\Nat$.  (T2): for an
instance $c=d$ of \rulename{I0} or \rulename{IS} the source $c$ is
iterator-headed and congruence in the template $W(v)$ gives $W(c)=W(d)$;
and from $W(s)=W(t)$ and $W(h(s))=W(r)$, \cref{guard:lem:wlift} gives
$W(h(s))=W(h(t))$, and symmetry and transitivity at the source $W(h(s))$
give $W(r)=W(h(t))$.  (T3): for a protected $p=q$,
\cref{guard:lem:safe}(v) with $a\defeq W(p)$ and $b\defeq W(q)$ gives
$(p=q)\dashv\vdash(W(p)=W(q))$.  (T4) is \cref{guard:lem:zero}.  Apply
\cref{guard:thm:public,guard:cor:biinterp}; the contradiction
$\varepsilon(0=1)$ is $\bot_W$.
\end{proof}

The compatibility of the two reductions is obtained here without
recovering symmetry at the source $0$: every equation that must later act
as a major premise is represented at stable iterator roots, the zero
computations are joined through $W(0)$, and the equalities of the native
derivation are moved by \cref{guard:lem:wlift}.  \Cref{guard:prop:inert}
does not apply because no public equation is zero-left.  No statement is
made about an auxiliary hypothesis $0=t$ in $\mathsf H$.

\begin{remark}[The executable extension and the form in Lean]
\label{guard:rem:lean-public}
The compact checkers of \pkg{protected\_core}, and those of
\pkg{guarded\_source}, \pkg{guarded\_orbit},
\pkg{deterministic\_auxiliary} and \pkg{mirror\_guard}, accept derivations
of $\mathsf U^{-}$.  Since
$\mathsf H\subseteq\mathsf U^{-}\subseteq\Nat$ as sets of derivable
sequents, by inclusion and \cref{guard:lem:incl}, the public consequences
of $\mathsf U^{-}$ on $\mathcal P_W$ are those of $\mathsf H$
(\cref{guard:thm:protected}(iii)), and likewise on $\mathcal P_m$ by
\cref{guard:thm:mirror}(iii).  \pkg{protected\_core} and \pkg{mirror\_guard}
note that for this inclusion the ordinary reading suffices and the mixed
reading is not needed.

In Lean the public languages of this section are formalized in a common
form (\Lean{NativePublic}).  The calculus \Lean{PGPrf} is the kernel with
the root-determined introduction and with transport at iterator-headed
sources only.  A gating translation \Lean{pub} replaces every equation
$s=t$ by the closure of its gate, $\Q{w}G_w(s=t)$, so that every atomic
matrix is guarded (\Lean{NFm.pub\_guarded}) and the two sides of an opened
gate are iterators, and
\Lean{pgprf\_pub\_iff}: $\Lean{PGPrf}\ (\Gamma.\mathrm{map}\ \mathit{pub})\
(\mathit{pub}\ A)\leftrightarrow\Lean{NPrf}\ \Gamma\ A$.  The theorem is
proved for any gating of equations (\Lean{Gating}); the mirror gate
$s=\It{w}{s}{u}{t}$ is a second instance (\Lean{Gating.mirror},
\Lean{NFm.pub\_mirrored}), with iterator-headed sources, which is stronger
than the exclusion of the literal $0$ in \cref{guard:thm:mirror}.  The
variant \Lean{NFm.pubD} uses the paired binder with an unused witness for
every closure and reaches the grammar of the packages
(\Lean{pgprf\_pubD\_iff}, \Lean{NFm.pubD\_noQ},
\Lean{NFm.pubD\_guarded}).  The proofs use neither \Lean{Classical.choice}
nor the bi-interpretation.  The differences from the packages are: an
equation at the root is gated, not left bare or wrapped as $W(s)=W(t)$;
clause contexts and the auxiliary formulas of \cref{guard:sec:aux} and
\cref{guard:sec:orbit} are not formalized; and the certificate formats are
not modelled.
\end{remark}

% ---------------------------------------------------------------------
\section{Guards on the full native language}
\label{guard:sec:full}

This section draws on \pkg{diagonal\_guard} for the criterion, the diagonal
guard and the lemmas on cases and implication, and on \pkg{mirror\_core}
for the mirror selector, its equational decoder, the repair of
substitution and the six-rule form.  \pkg{mirror\_core} takes the repair of
substitution from \pkg{reflexive\_base} and the mirrored shape from
\pkg{mirror\_guard}.

In this section the language is the full native language, the contexts are
lists of native formulas, and the readings of $\Q{}$ and $\D{}{}$ are the
ordinary ones.  The progress mode of the single introduction is used at the
guarded shape and nowhere else, so the atomic matrices that are not guarded
keep their introduction by cases, and \cref{guard:prop:strict} does not
arise.  The equality rule is the rule \rulename{Tr} of the kernel.

\subsection{Selectors and the criterion}

\begin{definition}[Selector]
\label{guard:def:selector}
A \emph{selector} is a property $P$ of matrices $B(x,y)$ that holds only of
literal equations, is decided from the syntax of $B$, and is invariant under
renaming of bound variables and under injective renaming of the variables
other than $x$ and $y$.  Put $\eta^{P}_B(z)\defeq z$ if $P(B)$ and
$\eta^{P}_B(z)\defeq 0$ otherwise.  The calculus $\Nat_P$ has the rules
\rulename{Hyp}, \rulename{I0}, \rulename{IS}, \rulename{Tr}, \rulename{QE}
and \rulename{DE} of $\Nat$ (\cref{base:sec:kernel}) and the one
introduction
\[
 \infrule{\Gamma,\ B(t,\eta^{P}_B(z))\proves B(t,\Suc z)}
         {\Gamma\proves\D{x}{y}B(x,y)}\ \rulename{DI}_{P},
\]
with $z$ not free in $\Gamma$, $t$ and $\D{x}{y}B$.  It has seven rules.
The test $P$ is applied to the matrix before the witness is substituted.
\end{definition}

\begin{lemma}[Structure and equality in $\Nat_P$]
\label{guard:lem:np-struct}
$\Nat_P$ admits weakening, injective renaming of free variables and cut.
Reflexivity, symmetry, transitivity, congruence and transport in an equation
template are derivable.
\end{lemma}

\begin{proof}
The structural statements are proved by induction on derivations as in
\cref{guard:lem:struct}: an injective renaming preserves $P$ and its
negation, and weakening and cut do not change a matrix.  The equality rules
are \cref{kern:lem:local}(iii): their derivations use \rulename{I0} and
\rulename{Tr} only.
\end{proof}

Consider two conditions on a selector.
\begin{enumerate}[label=\textup{(S\arabic*)}]
\item If $P(B)$, then $\Gamma\proves B(t,0)$ for every $\Gamma$ and $t$, by
 \rulename{I0}, \rulename{IS} and \rulename{Tr}.
\item Every equation template $E(y)$ has a \emph{lift} $\Lambda_E(y)$: an
 equation with $P(\Lambda_E)$, as a matrix in $y$ with an unused witness
 coordinate, such that $\Lambda_E(\Suc v)\dashv\vdash E(\Suc v)$ for every
 term $v$, by \rulename{I0}, \rulename{IS} and \rulename{Tr}.
\end{enumerate}

\begin{lemma}[Cases from \textup{(S1)}]
\label{guard:lem:np-cases}
Let $P$ satisfy \textup{(S1)}.  In $\Nat_P$, for every native matrix $B$ and
every native formula $F(y)$, with $z$ fresh:
\begin{enumerate}[label=\textup{(\roman*)}]
\item if $\Gamma,B(t,0)\proves B(t,\Suc z)$, then
 $\Gamma\proves\D{x}{y}B$;
\item if $\Gamma,F(0)\proves F(\Suc z)$, then $\Gamma\proves\Q{y}F(y)$;
\item if $\Gamma\proves F(0)$ and $\Gamma,F(0)\proves F(\Suc z)$, then
 $\Gamma\proves F(t)$ for every term $t$.
\end{enumerate}
\end{lemma}

\begin{proof}
(i) If $P(B)$ fails, this is $\rulename{DI}_P$.  If $P(B)$ holds, cut the
derivation of $B(t,0)$ given by (S1) into the premise, weaken by $B(t,z)$,
and apply $\rulename{DI}_P$ in progress mode.  (ii) and (iii) By (i) the
rule \rulename{DI} is admissible in $\Nat_P$, and the introduction of
$\Q{}$ by cases and numerical cases follow from \rulename{DI}, \rulename{DE},
\rulename{Hyp} and \rulename{QE} as in $\Nat$ (\cref{kern:lem:local}(i)).
\end{proof}

No induction for an equation is used in the lemma.  For equations
$U\equiv(a=b)$ and $V\equiv(c=d)$ and a fresh $w$ put
\[
 (U\Rightarrow_{=}V)\defeq\Q{w}\bigl(\sel(w,a,c)=\sel(w,b,d)\bigr).
\]
This is the native implication $U\Rightarrow_*V$ of \cref{base:def:macros}
at two equations.  Its rules are proved for $\Nat$ in
\cref{base:sec:structural}; the next lemma derives them in $\Nat_P$, where
\rulename{QI} is not yet available.

\begin{lemma}[Implication between equations, and cases on a term of the
context]
\label{guard:lem:np-imp}
Let $P$ satisfy \textup{(S1)}.  In $\Nat_P$:
\begin{enumerate}[label=\textup{(\roman*)}]
\item if $\Gamma,U\proves V$, then $\Gamma\proves U\Rightarrow_{=}V$; and
 $\Gamma\proves U\Rightarrow_{=}V$ and $\Gamma\proves U$ imply
 $\Gamma\proves V$;
\item for an equation $E$ and a term $r$, with $v$ not free in $\Gamma$,
 $r$ and $E$: if $\Gamma,r=0\proves E$ and $\Gamma,r=\Suc v\proves E$, then
 $\Gamma\proves E$.  The context $\Gamma$ may contain $r$ and its
 variables.
\end{enumerate}
\end{lemma}

\begin{proof}
(i) Let $J(w)$ be the matrix of $U\Rightarrow_{=}V$.  By the selector
equations $J(0)\dashv\vdash U$ and $J(\Suc z)\dashv\vdash V$.  For the
introduction, over $\Gamma,J(0)$ convert $J(0)$ to $U$, cut into the given
derivation, convert $V$ to $J(\Suc z)$, and apply
\cref{guard:lem:np-cases}(ii).  For the elimination, convert $U$ to $J(0)$,
apply \rulename{QE} at $1$ and convert $J(1)$ to $V$.

(ii) For a new variable $y$ put $F(y)\defeq\bigl((r=y)\Rightarrow_{=}E\bigr)$.
By (i) the two hypotheses give $\Gamma\proves F(0)$ and
$\Gamma\proves F(\Suc v)$.  \Cref{guard:lem:np-cases}(iii) at the target $r$
gives $(r=r)\Rightarrow_{=}E$, and reflexivity and (i) give $E$.  The
variable of the case distinction is the new $y$, not a variable of
$\Gamma$.
\end{proof}

\begin{theorem}[The selector criterion]
\label{guard:thm:selector}
Let $P$ be a selector with \textup{(S1)} and \textup{(S2)}.  Then the rule
\rulename{QI} of $\Nat$ is admissible in $\Nat_P$, and for every context
$\Gamma$ and every native formula $A$,
\[
 \Gamma\proves_{\Nat_P}A\iff\Gamma\proves_{\Nat}A .
\]
Both transformations are effective and change no term, formula or
assumption.  Hence $\Nat_P$ is bi-interpretable and equiconsistent with
$\iPRA$ by the interpretations of $\Nat$ (\cref{base:thm:bridge}).
\end{theorem}

\begin{proof}
From $\Nat_P$ to $\Nat$: six rules are shared.  An inference
$\rulename{DI}_P$ at a matrix without $P$ is \rulename{DI}.  At a matrix
with $P$ the matrix is a literal equation, and the inference is derived as
in \cref{guard:lem:incl}: \rulename{QI}, \rulename{QE} and \rulename{DI}.

From $\Nat$ to $\Nat_P$: \rulename{DI} is \cref{guard:lem:np-cases}(i), and
it remains to derive \rulename{QI}.  Let $\Gamma,E(z)\proves E(\Suc z)$
with $z$ fresh, let $\Lambda$ be the lift of $E$, and put
$\Delta\defeq\Gamma,E(0)$.

\emph{Decoding the lift.}  $\Delta,\Lambda(z)\proves E(z)$, by
\cref{guard:lem:np-imp}(ii) for the term $z$.  In the branch $z=0$,
transport of the hypothesis $E(0)$ along $0=z$ gives $E(z)$.  In the branch
$z=\Suc w$, transport of $\Lambda(z)$ gives $\Lambda(\Suc w)$, (S2) gives
$E(\Suc w)$, and transport along $\Suc w=z$ gives $E(z)$.

\emph{Progress of the lift.}  Cut $E(z)$ into the step premise:
$\Delta,\Lambda(z)\proves E(\Suc z)$, and by (S2)
$\Delta,\Lambda(z)\proves\Lambda(\Suc z)$.  The variable $z$ is not free in
$\Delta$.

\emph{Conclusion.}  Since $P(\Lambda)$, the rule $\rulename{DI}_P$ in
progress mode with the witness $0$ gives $\Delta\proves\D{x}{y}\Lambda(y)$,
and \rulename{DE} and \rulename{Hyp} give $\Delta\proves\Q{y}\Lambda(y)$.
With $\Lambda(0)$ from (S1), \rulename{QE} at $\Suc z'$ gives
$\Lambda(\Suc z')$, and (S2) gives $E(\Suc z')$.  Thus
$\Gamma,E(0)\proves E(\Suc z')$, and \cref{guard:lem:np-cases}(ii) gives
$\Gamma\proves\Q{y}E(y)$.

The base $E(0)$ is retained and discharged at the end: the conditional
closure is never identified with unconditional progress, and the argument
does not need $E(0)$ to be derivable.
\end{proof}

\begin{proposition}[Substitution in $\Nat_P$]
\label{guard:prop:np-subst}
\begin{enumerate}[label=\textup{(\roman*)}]
\item If $P$ is preserved and reflected by substitution, then substitution
 is admissible in $\Nat_P$ by an induction that keeps the rule and the mode
 of every inference.
\item If $P$ is preserved by substitution and satisfies \textup{(S1)}, then
 substitution is admissible in $\Nat_P$.
\end{enumerate}
\end{proposition}

\begin{proof}
Induct on the derivation, renaming each eigenvariable away from the values
of the substitution.  All rules other than $\rulename{DI}_P$ commute with
substitution as in $\Nat$.  For $\rulename{DI}_P$ in progress mode,
preservation gives the same mode after substitution.  For $\rulename{DI}_P$
in case mode at $B$: in (i) the substituted matrix does not have $P$, by
reflection, and the same inference applies; in (ii) the induction hypothesis
gives the substituted premise by cases, and \cref{guard:lem:np-cases}(i)
derives the conclusion whether or not the substituted matrix has $P$.
\end{proof}

\begin{remark}[Versions of the proof]
\label{guard:rem:selector-versions}
The criterion is the \emph{guard criterion} of \pkg{diagonal\_guard}.  That
package derives \rulename{QI} with the implication
$\Lambda(y)\Rightarrow_{=}E(\Suc y)$, proved by cases on a new variable, and
needs the step premise at $0$ and at $\Suc v$; it substitutes into the
derivation of the premise, which is justified for the diagonal selector by
\cref{guard:prop:np-subst}(i) and, for a selector without the stability
property, by translating native derivations by induction on their height
and substituting into the native premise first.  \pkg{mirror\_core} decodes
the lift by two unconditional equations (\cref{guard:lem:decoder}) and uses
the step premise at its own eigenvariable.  The proof above takes the order
of \pkg{mirror\_core} and decodes with \cref{guard:lem:np-imp}(ii), which is
a lemma of \pkg{diagonal\_guard}; it needs neither a substitution nor a
special property of the lift.
\end{remark}

\subsection{The diagonal guard}

\begin{theorem}[The diagonal guard]
\label{guard:thm:diagonal}
Let $P$ hold of the matrices that are diagonal in $y$
(\cref{guard:def:shapes}), and write $\Nat_{\mathrm{diag}}$ for $\Nat_P$.
\begin{enumerate}[label=\textup{(\roman*)}]
\item $P$ is a selector that satisfies \textup{(S1)} and \textup{(S2)}, the
 latter with the diagonal lift $\Delta_yE$.
\item $\Gamma\proves_{\Nat_{\mathrm{diag}}}A$ if and only if
 $\Gamma\proves_{\Nat}A$, for all native $\Gamma$ and $A$.
\item $\Nat_{\mathrm{diag}}$ admits weakening, substitution and cut, and
 the substitution keeps the rule and the mode of every inference.
\end{enumerate}
\end{theorem}

\begin{proof}
(i) (S1): at $y=0$ both sides of a diagonal matrix have the form
$\It{0}{0}{u}{h'}$ and are equal to $0$ by \rulename{I0}.  (S2): the lift
$\Delta_yE$ is diagonal in $y$, and at $y=\Suc v$ its two sides are
$\It{\Suc v}{\Suc v}{u}{s(\Suc v)}$ and $\It{\Suc v}{\Suc v}{u}{t(\Suc v)}$,
which are equal to $s(\Suc v)$ and $t(\Suc v)$ by \rulename{IS}, the updates
being constant.  No identity between arbitrary iterator bodies is asserted.
(ii) is \cref{guard:thm:selector}.  (iii) By
\cref{guard:prop:subst-shapes}(iii) the diagonal shape is preserved and
reflected, so \cref{guard:prop:np-subst}(i) and \cref{guard:lem:np-struct}
apply.
\end{proof}

In Lean the calculus is \Lean{GPrf} (\Lean{NativeDiagonal}), with
\Lean{gprf\_iff\_nprf}, \Lean{gprf\_iff\_iprf} and substitution by direct
induction (\Lean{GPrf.subst}); the equivalence uses neither
\Lean{Classical.choice} nor the bi-interpretation.  The formalization
recovers progress by an equational decoder in place of the implication of
the package: with $E(y)\equiv(s(y)=t(y))$, numerical cases prove the
unconditional equation
\[
 \sel\bigl(y,\ s(0),\ \It{y}{y}{u}{s(y)}\bigr)=s(y)
\]
and the same for $t$, and from $E(0)$ and $\Delta_vE$ they give
$s(v)=\sel(v,s(0),\It{v}{v}{u}{s(v)})=\sel(v,t(0),\It{v}{v}{u}{t(v)})=t(v)$.

The diagonal guard keeps seven rules, the full language and the readings of
both binders.  The witness-selected calculus $\mathsf W$
(\cref{kern:def:W,kern:thm:witness}) also has seven rules on the full
language, and its mode can change when a substitution turns a witness into
the literal $0$ (\cref{kern:rem:wsubst}); the calculus with orbit
elimination (\cref{kern:sec:R}) changes the reading of $\Q{}$ at equations
and its elimination; and the guarded core restricts the public language.  The
diagonal guard does none of these.

\subsection{The mirror selector and six rules}

\begin{theorem}[The mirror selector]
\label{guard:thm:mirror-selector}
Let $\mathrm{Mir}$ hold of the matrices that are mirrored in $y$, and write
$\Nat_{\mathrm{mir}}$ for $\Nat_{\mathrm{Mir}}$.
\begin{enumerate}[label=\textup{(\roman*)}]
\item $\mathrm{Mir}$ is a selector that satisfies \textup{(S1)} and
 \textup{(S2)}, the latter with the lift
 $H_yE(y)=\bigl(s(y)=\sel(y,s(y),r(y))\bigr)$ of $E(y)\equiv(s(y)=r(y))$.
\item $\Gamma\proves_{\Nat_{\mathrm{mir}}}A$ if and only if
 $\Gamma\proves_{\Nat}A$.
\item Substitution is admissible in $\Nat_{\mathrm{mir}}$.  It does not
 keep the mode of every inference.
\end{enumerate}
\end{theorem}

\begin{proof}
(i) (S1): the instance at zero is $a'=\It{0}{a'}{u}{h'}$, the reverse of an
instance of \rulename{I0}.  (S2): $\sel(\Suc v,s,r)=r$ by \rulename{IS}.
(ii) is \cref{guard:thm:selector}.  (iii) The shape is preserved
(\cref{guard:prop:subst-shapes}(ii)), so \cref{guard:prop:np-subst}(ii)
applies.  It is not reflected: the matrix $p=\It{y}{q}{u}{u}$ is not
mirrored for distinct variables $p,q$ and is mirrored after both are
replaced by one variable, so an introduction by cases at this matrix
becomes, under that substitution, an instance of
\cref{guard:lem:np-cases}(i), which cuts the provable base and weakens to
progress.
\end{proof}

The selector is tested on the matrix before the witness is substituted: the
matrix $x=\It{y}{0}{u}{0}$ is not mirrored, although its instance at the
witness $0$ is $0=\It{y}{0}{u}{0}$.

\begin{lemma}[The equational decoder of the mirror lift]
\label{guard:lem:decoder}
Let $E(y)\equiv(s(y)=r(y))$ and $\Lambda(y)\equiv
\bigl(s(y)=\sel(y,s(y),r(y))\bigr)$.  In $\Nat_{\mathrm{mir}}$, for every
term $v$, $\Gamma\proves E(0)$ and $\Gamma\proves\Lambda(v)$ imply
$\Gamma\proves E(v)$, by numerical cases for equations and equality alone.
\end{lemma}

\begin{proof}
The equations
\begin{align*}
 \sel\bigl(v,\ s(0),\ s(v)\bigr)&=s(v),\\
 \sel\bigl(v,\ r(0),\ \sel(v,s(v),r(v))\bigr)&=r(v)
\end{align*}
are derivable over the empty context by \cref{guard:lem:np-cases}(iii): at
$0$ both left sides are equal to their seed by \rulename{I0}, and at a
successor the first uses one application of \rulename{IS} and the second
two.  They are proved before any assumption about $v$ is retained, and then
weakened.  From $\Lambda(v)$, congruence in the template $\sel(v,s(0),-)$
gives $\sel(v,s(0),s(v))=\sel(v,s(0),\sel(v,s(v),r(v)))$, and the base
$s(0)=r(0)$ changes the seed on the right.  Hence
\[
 s(v)=\sel\bigl(v,s(0),s(v)\bigr)
 =\sel\bigl(v,s(0),\sel(v,s(v),r(v))\bigr)
 =\sel\bigl(v,r(0),\sel(v,s(v),r(v))\bigr)=r(v).
\]
\end{proof}

The lemma replaces the first step of the proof of
\cref{guard:thm:selector} for the mirror selector, without an implication
between equations.

\begin{definition}[The six-rule mirror calculus]
\label{guard:def:m6}
The calculus $\mathsf M_6$ is the calculus $\mathsf E_2$ of
\cref{comp:def:e3} with its introduction \rulename{WI} replaced by the
mirror-selected one: its rules are \rulename{Hyp}, \rulename{Refl}, the two
cases $\rulename{Eval}_0$ and $\rulename{Eval}_{\Suc}$ of \rulename{Eval},
\rulename{QE}, \rulename{DE} and
\[
 \infrule{\Gamma,\ B(t,\eta^{\mathrm{Mir}}_B(z))\proves B(t,\Suc z)}
         {\Gamma\proves\D{x}{y}B(x,y)}\ \rulename{DI}_{\mathrm{Mir}} .
\]
It has six rule families and seven inference constructors, counted as in
\cref{comp:rem:counting}.  There is no computation axiom, no transport rule,
no structural rule and no rule of induction.
\end{definition}

\begin{theorem}[Six rules with the mirror selector]
\label{guard:thm:m6}
For every context $\Gamma$ and every native formula $A$,
\[
 \Gamma\proves_{\mathsf M_6}A\iff\Gamma\proves_{\Nat_{\mathrm{mir}}}A
 \iff\Gamma\proves_{\Nat}A .
\]
$\mathsf M_6$ admits weakening, injective renaming, cut and substitution,
and is bi-interpretable and equiconsistent with $\iPRA$ by the
interpretations of $\Nat$.
\end{theorem}

\begin{proof}
The derivations of symmetry and transitivity, of \rulename{I0} and
\rulename{IS}, and of \rulename{Tr} in $\mathsf E_2$
(\cref{comp:lem:evalalg,comp:lem:evalcomp,comp:cor:evaltr}) use
\rulename{Refl} and the two cases of \rulename{Eval} only, so they are
derivations of $\mathsf M_6$; the other rules of $\Nat_{\mathrm{mir}}$ are
rules of $\mathsf M_6$.  Conversely \rulename{Refl} and both cases of
\rulename{Eval} are derived from \rulename{I0}, \rulename{IS} and
\rulename{Tr} as in the proof of \cref{comp:thm:eval}, hence in
$\Nat_{\mathrm{mir}}$.  The second equivalence is
\cref{guard:thm:mirror-selector}.  Weakening, injective renaming and cut
are proved for $\mathsf M_6$ as for $\mathsf E_3$ in the proof of \cref{comp:thm:eval}, with the
observation that an injective renaming preserves and reflects
$\mathrm{Mir}$.  Substitution is proved as in
\cref{guard:prop:np-subst}(ii): \rulename{Eval} commutes with substitution
(\cref{comp:prop:phisubst}), and the derived introduction by cases of
\cref{guard:lem:np-cases}(i) is available in $\mathsf M_6$ by the first
equivalence.
\end{proof}

\begin{remark}[The derivations of \pkg{mirror\_core}]
\label{guard:rem:m6-versions}
\pkg{mirror\_core} derives the equality rules of $\mathsf M_6$ directly.
It joins $s=t$ and $s=r$ to $r=t$ through $\It{0}{s}{u}{u}$ by two
instances of $\rulename{Eval}_0$, which is the identity frame of
\cref{comp:lem:frame}, and it derives \rulename{Tr} from the two iterators
$\It{0}{s}{v}{h(v)}$ and $\It{1}{s}{v}{h(v)}$ by the composition described
in the remark after \cref{comp:cor:evaltr}.
\end{remark}

In Lean the seven-rule calculus $\Nat_{\mathrm{mir}}$ is \Lean{MiPrf}
(\Lean{miprf\_iff\_nprf}) and the six-rule calculus is \Lean{M6Prf}
(\Lean{m6prf\_iff\_nprf}), in the module \Lean{NativeMirror}.  The
structural rules proved there are weakening, injective renaming
(\Lean{MiPrf.ren}) and cut, which suffice for weakening below a binder and
for the equivalence; the progress of an arbitrary equation uses the lift
$s=\sel(y,s,t)$.  Relative to the six-rule calculus with the diagonal
selector (\Lean{SPrf}), the mirror selector shows one iterator where the
diagonal shows two, and it loses the reflection of the shape.

% ---------------------------------------------------------------------
\section{Comparison of the guards}
\label{guard:sec:compare}

This section compares the guards of \pkg{guarded\_core},
\pkg{protected\_core}, \pkg{mirror\_guard}, \pkg{mirror\_core} and
\pkg{diagonal\_guard}; \cref{guard:tab:guards} restates results of the
preceding sections.

\begin{table}[ht]
\centering
\small
\begin{tabular}{@{}p{0.13\linewidth}p{0.2\linewidth}p{0.17\linewidth}p{0.19\linewidth}p{0.2\linewidth}@{}}
\toprule
Guard & Atomic matrix & Lift of $s=t$ & Under substitution & Used in\\
\midrule
zero-seeded & $\It{y}{0}{u}{h}=\It{y}{0}{u'}{k}$ &
 $\sel(y,0,s)=\sel(y,0,t)$ &
 preserved, not reflected &
 $\mathcal P_0$: $\mathsf G$, $\mathsf U^{-}$,
 $\mathsf G^{-}_{\mathrm{orb}}$\\
wrapped & zero-seeded; root atoms protected &
 $W(s)=W(t)$ at the root &
 preserved; sources closed &
 $\mathcal P_W$: $\mathsf H$\\
mirror & $a=\It{y}{a}{u}{h}$ &
 $s=\sel(y,s,t)$ &
 preserved, not reflected &
 $\mathcal P_m$: $\mathsf M$; full language:
 $\Nat_{\mathrm{mir}}$, $\mathsf M_6$\\
diagonal & $\It{y}{y}{u}{h}=\It{y}{y}{u'}{k}$ &
 $\It{y}{y}{u}{s}=\It{y}{y}{u}{t}$ &
 preserved and reflected &
 full language: $\Nat_{\mathrm{diag}}$\\
\bottomrule
\end{tabular}
\caption{The four guards.}
\label{guard:tab:guards}
\end{table}

\emph{The base.}  Every guard has an instance at zero that is derivable by
\rulename{I0} and equality at iterator-headed sources
(\cref{guard:lem:zero}).  For the mirror guard it is the reverse of one
computation equation.  For the zero-seeded and the diagonal guard two
computations are joined; under a source restriction the join passes through
$W(0)$ (\cref{guard:lem:safe}(iv)).

\emph{The left side.}  The zero-seeded, wrapped and diagonal shapes have an
iterator on the left.  With the sources $s\not\equiv 0$ and the uncrossed
transport this makes a zero-left hypothesis inert
(\cref{guard:prop:inert}): the zero-seeded guard then needs an unguarded
matrix in the proof language (\cref{guard:thm:aux,guard:thm:orbit}), or
public atoms that are never zero-left (\cref{guard:thm:protected}).  The
mirror guard leaves the left side free and is complete in a strictly
guarded language (\cref{guard:thm:mirror}).

\emph{Sources.}  The guarded core uses all sources.  The mirror guard, and
the zero-seeded guard with auxiliary or orbit formulas, use the sources
$s\not\equiv 0$, a class that is not closed under substitution; their
proofs avoid substitution until equality is unrestricted.  The protected
core uses iterator-headed sources, a class closed under substitution, and
pays with a wrapper on every public atom and with a larger validator of the
public grammar (\cref{guard:sec:checkers}).

\emph{Substitution of the shape.}  In a public language the mode is
selected by the root and the shape is a condition on the language, so
preservation is enough.  As a selector on the full language the shape must
not be created by substitution if the mode of an inference is to be stable:
the diagonal shape is the only one of the four with this property
(\cref{guard:prop:subst-shapes}), and for it substitution is a
rule-preserving induction (\cref{guard:thm:diagonal}).  The mirror selector
is closed under substitution with a repair
(\cref{guard:thm:mirror-selector}).  The zero-seeded shape also satisfies
\textup{(S1)} and \textup{(S2)}, with the gate $G_yE$ as the lift, so by
\cref{guard:thm:selector,guard:prop:np-subst} it is a selector on the full
language with the same repair; \pkg{diagonal\_guard} records the matrix
$\It{y}{p}{u}{0}=\It{y}{0}{u}{0}$ as the witness that its mode is not
stable.

\emph{Cost of recognition.}  The zero-seeded and the diagonal test inspect
fixed positions of the two sides.  The mirror test compares two term trees,
which is linear in their size; \pkg{mirror\_core} notes that its smaller
checker source does not imply faster checking or shorter proofs.

% ---------------------------------------------------------------------
\section{Checkers and their sizes}
\label{guard:sec:checkers}

This section records the programs of \pkg{guarded\_core},
\pkg{guarded\_source}, \pkg{guarded\_orbit},
\pkg{deterministic\_auxiliary}, \pkg{mirror\_guard}, \pkg{protected\_core},
\pkg{mirror\_core} and \pkg{diagonal\_guard}.  Sizes are measured by the
protocol of \cref{base:sec:measure}.

\paragraph{The compact checkers.}
\pkg{guarded\_core}, \pkg{guarded\_source}, \pkg{guarded\_orbit},
\pkg{deterministic\_auxiliary}, \pkg{mirror\_guard} and
\pkg{protected\_core} modify the arity-coded checkers of the earlier
packages; the codec is inherited and is not changed.  A formula array of
length one is $\D{}{}$, of length two an equation, of length three an
administrative $\Q{}$; a term array of length one is a successor and of
length three an iterator; the other lengths denote zero or truth.  The
root-determined introduction replaces the comparison of the parity of the
matrix with a field of the certificate, \texttt{len(a) \% 2 == u}, by the
parity alone, \texttt{not len(a) \% 2}: among the meaningful roots only the
equation has even length, so an equation retains its instance at the
eigenvariable and a compound matrix its instance at zero.
\pkg{guarded\_source} writes the test as \texttt{len(a) == 2}, and
\pkg{deterministic\_auxiliary} has an optional spelling
\texttt{\textasciitilde len(a) \% 2} that saves three bytes.  The byte
profile of \pkg{guarded\_core} also moves the check of the major premise
into the expression that supplies the hypothesis, which preserves every
successful run.  Source selection is the test
\texttt{m == ([s, u] if s else [t] * 3)}: the canonical literal $0$ is the
integer $0$, which is false, and selects the $\Q{}$ major, and every other
source selects the equation.  Computation packets that cannot have a
literal $0$ as initializer or as predecessor are obtained from the
\emph{certified zero} $c_0\defeq\It{2}{1}{u}{0}$, with $c_0=0$, by a
transport at the iterator-headed source $c_0$, so that no transport at the
literal $0$ is reintroduced.  The byte and AST profiles use lexical scopes
and the token profiles a global supply of fresh names; their certificates
are not interchangeable.

\paragraph{The contract of the compact checkers.}
A compact checker accepts derivations of $\mathsf U$
(\pkg{guarded\_core}) or of $\mathsf U^{-}$ (the other five) in the
administrative syntax.  Each package states that, for a closed public
formula $A$ in canonical encoding and in the ideal unbounded execution, an
accepting certificate exists if and only if $A$ is derivable in its
calculus.  The soundness half rests on the correctness of the inherited
codec (opening, freshness and the readings of the surplus arities), which
these packages cite and do not reprove, together with
\cref{guard:cor:admin,guard:thm:aux,guard:thm:orbit,guard:thm:mirror,guard:thm:protected}
and \cref{guard:rem:lean-public}; the completeness half encodes a
derivation with fresh integer names.  A compact checker does not test that
the goal is public, and it accepts some goals that are not.  Each package
therefore supplies a \emph{strict} checker, which validates the closed
public goal and then runs the same kernel, and counts it in full.  The
strict checker of \pkg{mirror\_guard} compares the two occurrences of the
seed after raising the left side across the binder of the iterator; a
comparison of the raw arrays would reject matrices with bound variables.
One inherited negative control that changes only the retired mode field is
accepted by the new checkers; an old certificate that introduces an
unguarded equation by cases is rejected by them, which is
\cref{guard:prop:strict} and the reason why the raw certificate languages
of the old and the new checkers are not claimed to coincide.

\paragraph{The typed checkers.}
\pkg{mirror\_core} and \pkg{diagonal\_guard} measure checkers with tagged
terms, formulas and proof trees in de Bruijn form, which validate every
term, formula, scope and proof node of the full native language.  The
diagonal checker tests, after opening the witness, whether both sides begin
with an iterator whose counter and seed are the bound index $0$; by
\cref{guard:prop:subst-shapes}(iii) the test after opening is equivalent to
the test on the matrix.  The mirror checker tests the matrix before the
witness is substituted.  The two packages differ from their controls in the
guard of the introduction only.  The translator of \pkg{diagonal\_guard}
implements both directions of \cref{guard:thm:diagonal}(ii); three
translated certificates grow from $485$, $885$ and $1{,}965$ bytes to
$11{,}506$, $14{,}218$ and $20{,}283$ bytes.  \pkg{mirror\_core} supplies
the translation from $\mathsf M_6$ to the kernel, the derived introduction
by cases and the repaired substitution as programs; the translation of
\rulename{QI} into $\mathsf M_6$ is given by the written proof only.

\paragraph{Recorded runs.}
The packages record the following executions of their programs on positive
and negative certificates, in ordinary and optimized Python, all with the
expected outcome: $800$ (\pkg{guarded\_core}), $150$
(\pkg{guarded\_source}), $224$ (\pkg{guarded\_orbit}), $242$
(\pkg{deterministic\_auxiliary}), $480$ (\pkg{mirror\_guard}), $608$
(\pkg{protected\_core}); for \pkg{mirror\_core}, $92$ command-line checks of
$23$ fixtures, $36$ replays in the native checker and $72$ replays of
numerical substitutions; for \pkg{diagonal\_guard}, $192$ processes on $32$
fixtures, $93$ validations of translated certificates and $240$ finite
checks of the preservation and reflection of the guard.  The runs test the
programs and the certificates; the theorems of this chapter are the written
proofs.

\begin{longtable}{@{}llrrrl@{}}
\caption{Programs and sizes as the packages state them.  The last column is
the matched control with its tokens, AST nodes and normalized bytes.  The
smallest values of one package are attained by different programs.}
\label{guard:tab:sizes}\\
\toprule
Package & Program & Tokens & AST & Bytes & Matched control\\
\midrule
\endfirsthead
\toprule
Package & Program & Tokens & AST & Bytes & Matched control\\
\midrule
\endhead
\bottomrule
\endfoot
\pkg{guarded\_core} & direct & 265 & 308 & 528 & byte profile, 266/310/529\\
 & byte profile & 266 & 307 & 527 & byte profile, 266/310/529\\
 & token profile & 260 & 290 & 575 & token profile, 261/292/576\\
 & AST profile & 263 & 286 & 607 & AST profile, 263/289/613\\
 & readable & 292 & 336 & 1275 & none\\
 & strict & 491 & 568 & 1208 & old core with goal parser, 450/521/1109\\
\midrule
\pkg{guarded\_source} & token profile & 257 & 285 & 568 &
 \pkg{source\_dispatch}, 259/289/572\\
 & AST profile & 261 & 279 & 603 & \pkg{source\_dispatch}, 263/283/607\\
 & byte profile & 264 & 301 & 522 & \pkg{source\_dispatch}, 266/305/526\\
 & strict & 490 & 561 & 1202 & \pkg{guarded\_core} strict, 491/568/1208\\
\midrule
\pkg{guarded\_orbit} & token profile & 258 & 287 & 571 &
 \pkg{source\_dispatch}, 259/289/572\\
 & AST profile & 262 & 281 & 606 & \pkg{source\_dispatch}, 263/283/607\\
 & byte profile & 265 & 303 & 525 & \pkg{source\_dispatch}, 266/305/526\\
 & strict & 491 & 563 & 1205 & guarded 491/568/1208; scalar 450/521/1109\\
\midrule
\pkg{deterministic\_auxiliary} & direct & 265 & 303 & 525 &
 \pkg{source\_dispatch}, 266/305/526\\
 & byte profile & 265 & 303 & 522 & \pkg{source\_dispatch}, 266/305/526\\
 & AST profile & 262 & 281 & 606 & \pkg{source\_dispatch}, 263/283/607\\
 & token profile & 258 & 287 & 571 & \pkg{source\_dispatch}, 259/289/572\\
 & strict & 491 & 563 & 1205 & guarded 491/568/1208; native 450/521/1109\\
\midrule
\pkg{mirror\_guard} & token profile & 258 & 287 & 571 &
 \pkg{source\_dispatch}, 259/289/572\\
 & AST profile & 262 & 281 & 606 & \pkg{source\_dispatch}, 263/283/607\\
 & byte profile & 265 & 303 & 525 & \pkg{source\_dispatch}, 266/305/526\\
 & readable & 279 & 319 & 1272 & none\\
 & strict & 578 & 661 & 1503 & strict derivative of the source checker,
 471/537/1127\\
\midrule
\pkg{protected\_core} & token profile & 258 & 287 & 571 &
 \pkg{source\_dispatch}, 259/289/572\\
 & AST profile & 262 & 281 & 606 & \pkg{source\_dispatch}, 263/283/607\\
 & byte profile & 265 & 303 & 525 & \pkg{source\_dispatch}, 266/305/526\\
 & readable & 284 & 326 & 1301 & none\\
 & strict & 505 & 582 & 1251 & \pkg{guarded\_core} strict, 491/568/1208\\
\midrule
\pkg{mirror\_core} & mirror selector, six rules & 1012 & 1082 & 3109 &
 diagonal selector, six rules, 1022/1091/3125\\
\midrule
\pkg{diagonal\_guard} & diagonal, seven rules & 1069 & 1124 & 3244 &
 native, eight rules, 1101/1160/3382;\\
 & & & & & witness-selected, seven rules, 1041/1099/3205\\
\end{longtable}

The gains recorded in \cref{guard:tab:sizes} are small.  The compact
profiles of \pkg{guarded\_orbit},
\pkg{deterministic\_auxiliary}, \pkg{mirror\_guard} and
\pkg{protected\_core} have the same three sizes, since they share the two
dispatch tests and differ in the calculus that justifies them and in the
strict validator.  The strict checkers do not improve on the unrestricted
control of $1109$ bytes; those of \pkg{mirror\_guard} and
\pkg{protected\_core} are larger than their matched controls.  The gzip
size of the byte profile of \pkg{guarded\_orbit},
\pkg{deterministic\_auxiliary} (direct), \pkg{mirror\_guard} and
\pkg{protected\_core} is $331$ bytes against $330$ for the control.  The
diagonal checker is $39$ bytes larger than the witness-selected control; its
advantage is the stability of the mode under substitution.  The typed
checkers are not smaller than the compact ones, which have a different
input contract.

% ---------------------------------------------------------------------
\section{Further questions}
\label{guard:sec:questions}

This section collects the open questions and the unproved or negative
statements of \pkg{guarded\_core}, \pkg{guarded\_source},
\pkg{guarded\_orbit}, \pkg{deterministic\_auxiliary}, \pkg{mirror\_guard},
\pkg{protected\_core}, \pkg{mirror\_core} and \pkg{diagonal\_guard}.

\begin{question}[Minimal proof language]
\label{guard:q:minimal}
For the zero-seeded public language $\mathcal P_0$ with the sources
$s\not\equiv 0$, \cref{guard:cor:obstruction} shows that some unguarded
matrix is needed and \cref{guard:thm:orbit} shows that the orbit matrices
suffice.  Which substitution-stable families of auxiliary equation
invariants suffice for the repair of equality and for the arithmetic of the
bridge?  Is there a comparably short checker that enforces the smaller
intermediate grammar, with a simpler soundness proof than that of the full
arity codec?  (\pkg{guarded\_source}, \pkg{guarded\_orbit},
\pkg{deterministic\_auxiliary}.)
\end{question}

\begin{question}[The smallest stable guard]
\label{guard:q:stable}
Is there a syntactic class of equations that is cheaper to recognize than
the diagonal shape and has provable instances at zero, a lift of every
equation, and preservation and reflection under substitution?  What is the
smallest family of base recognizers that supports an equational decoder for
every equation?  The mirror and the diagonal shape are two instances;
neither gives a lower bound.  (\pkg{diagonal\_guard}, \pkg{mirror\_core}.)
\end{question}

\begin{question}[A public grammar with a small validator]
\label{guard:q:validator}
Is there a substitution-stable public grammar that makes the root-determined
introduction complete, keeps both composites, and has a strict checker no
larger than the unrestricted control of $1109$ bytes under the same
complete-input contract?  (\pkg{guarded\_orbit},
\pkg{deterministic\_auxiliary}; \pkg{guarded\_source} asks whether the
validator can share more of the kernel.)
\end{question}

\begin{question}[Zero-source symmetry]
\label{guard:q:symmetry}
Can the derivation of $t=0$ from $0=t$ with the uncrossed rule
(\cref{guard:thm:dependency}) be made shorter than the one through the
singleton test and the filter?  A shorter uniform derivation would shorten
certificates without changing a checker.  (\pkg{guarded\_source},
\pkg{deterministic\_auxiliary}.)
\end{question}

\begin{question}[Sizes of translated derivations]
\label{guard:q:proofsize}
Give explicit bounds for the expansion caused by the join at zero, the
padding, the implication between equations and the lift of an equation, and
compare proof trees with a checked format with sharing.  No polynomial bound
on the size of a translated certificate is claimed.
(\pkg{guarded\_orbit}, \pkg{diagonal\_guard}.)
\end{question}

The following items are stated by the packages without proof, or record what
was not achieved.

\begin{itemize}
\item \emph{No minimality.}  No package claims a least number of rules, a
 least checker, or a lower bound for a class of checkers; the observed sizes
 are upper bounds, and the three smallest coordinates of a package belong
 to different programs.  Whether one checker improves all three coordinates
 of its controls at once is open (\pkg{guarded\_source}), as is a lower
 bound for a precisely specified class of complete checkers
 (\pkg{mirror\_core}).
\item \emph{The codec.}  The theorems on the compact checkers are stated
 relative to the inherited arity codec, whose soundness is cited and not
 reproved in these packages, and neither the codec nor a certificate format
 is formalized in a proof assistant.
\item \emph{Administrative formulas.}  Completeness is not claimed for goals
 outside the public language, and the full language of $\mathsf U$ or
 $\mathsf U^{-}$ is not claimed to be bi-interpretable with $\iPRA$; by
 \cref{guard:prop:strict} the two sound readings of an unguarded atomic
 binder differ.  No impossibility is claimed for other interpretations of
 the administrative language (\pkg{guarded\_core}).  Whether the
 independent choice between cases and progress can be removed on every
 administrative formula with its ordinary reading and without a validator
 is asked by \pkg{guarded\_orbit}; the selectors of \cref{guard:sec:full}
 remove it on the full native language at the cost of a test of the shape.
\item \emph{Formalization.}  In Lean the dependency theorem
 (\Lean{ZKernel.trF}), the diagonal guard (\Lean{GPrf}), the mirror
 selector (\Lean{MiPrf}, \Lean{M6Prf}) and the public languages in the
 gated form (\Lean{PGPrf}) are proved.  Not formalized: the calculi with
 clause contexts; the auxiliary formulas of \cref{guard:sec:aux} and the
 orbit bootstrap of \cref{guard:sec:orbit}, including numerical cases at
 stable-left equations; the mirror calculus $\mathsf M$ with the sources
 $s\not\equiv 0$ in the form of \cref{guard:thm:mirror}; and the wrapped
 root atoms of $\mathsf H$.  \pkg{deterministic\_auxiliary} asks whether
 the auxiliary progress of an equation can be a separate proof judgment
 with a smaller derivation datatype, the target being the equality of the
 public consequence relations and not a round trip for auxiliary formulas;
 \pkg{mirror\_core} proposes its two unconditional selector equations
 (\cref{guard:lem:decoder}) as the route with the fewest dependencies.
 \pkg{mirror\_guard} observes that a
 formalization of $\mathsf M$ cannot start by restricting the transport
 rule of the kernel and using its substitution theorem
 (\cref{guard:rem:substwording}).
\item \emph{Programs that are not complete translators.}  The proof
 generators of \pkg{mirror\_guard} and \pkg{protected\_core} are collections
 of worked constructions; the normalizer of \pkg{protected\_core} rejects an
 unsupported equation leaf with the source $0$ instead of translating it;
 \pkg{mirror\_core} has no program for the translation of \rulename{QI}.
\item \emph{Selectors in the compact interface.}  Whether the mirror or the
 diagonal selector reduces the size of a compact checker with a complete
 validator of the goal is not answered (\pkg{mirror\_core},
 \pkg{diagonal\_guard}); nor is the combination of the diagonal introduction
 with the crossed one-term transport in a matched typed checker
 (\pkg{diagonal\_guard}), or a constant-time mirror test by sharing of
 syntax without a larger trusted reader (\pkg{mirror\_core}).
\item \emph{Other bases.}  \pkg{guarded\_orbit} asks whether a different
 numerical basis can remove a constructor or a binding operation without a
 richer arithmetic axiom or a more complex parser, and notes that the
 earlier register, tied-seed and zero-start results constrain the answer.
\end{itemize}
